% 本文件由 LaTeX 排版 Agent 根据有效 translation.json 自动生成。
% author=Augustine of Hippo; work_id=1703; title=独语录

\chapter{独语录}

% structure: p0001 source=h1 target=omitted reason=page-title-already-rendered-by-parent

\Emph{卷一}\newline{}\Emph{卷二}\par

\SourceRecord{Translated by C.C. Starbuck.；Nicene and Post-Nicene Fathers, First Series；Vol. 7.；Edited by Philip Schaff.；Buffalo, NY: Christian Literature Publishing Co.,；1888.；<http://www.newadvent.org/fathers/1703.htm>.}

% structure: p0001 source=h1 target=section reason=page-title

\section{独语录，第一卷}

当我长久独自思考许多不同的事情，并多日勤勉地探究关于我自己和我自己的至善，或者有什么恶是我要避免的：忽然有人对我说话——是我自己，或是别的某人，在我之内或之外，我不知道。因为这正是我主要努力要知道的。于是，有一个声音对我说，我们就称它为理性吧——\Person{理}：\Quote{请看，假设你已发现些什么，你将把它托付给谁，以便你能继续其他事情？}\Person{奥}：\Quote{记忆，无疑。}\Person{理}：\Quote{但记忆的力量如此强大，足以安全保存一切可能由思想产生的东西吗？}\Person{奥}：\Quote{很难，不，确实不能。}\Person{理}：\Quote{因此你必须写下来。但鉴于你的健康抗拒写作的辛劳，你又能做什么？这些事也不宜口述，因为它们只有完全独处时才能容许。}\Person{奥}：\Quote{对。因此我完全不知该说什么。}\Person{理}：\Quote{祈求天主赐你健康和助佑，使你更能达成你的愿望，并将这个祈求写下来，使你在脑力的产物中更加勇敢。然后，把你所发现的总结成几条简短的结论。现在不必操心去招徕众多读者；这些事若能为你本城的少数人听闻，也就足够了。}\par

2. 天主，宇宙的创造者，首先赐我正确地呼求祢；然后使我相称于被祢垂听；最后，俯允释放我。天主，通过祢，一切本不存在的事物得以趋向存在。天主，祢甚至连看似相互毁灭之物也保全不使其消亡。天主，祢从虚无中创造了这个世界，众目所见它极其美丽。天主，祢不造恶，却使恶不至最恶。天主，祢向少数投奔真实存在者显示恶乃是虚无。天主，通过祢，宇宙连其凶恶的一面也是完美的。天主，从祢而来，即使与祢最相异的事物也不造成不和谐，因为更差的事物也被纳入与更好事物同一计划中。天主，凡能爱之物，无论有意无意都爱祢。天主，万有都在祢内，然而任何受造物的卑劣在祢并非卑劣，其邪恶并非有害，其错误并非错误。天主，祢只愿纯洁者认识真理。天主，真理之父，智慧之父，真实而圆满生命之父，真福之父，美好与美丽之父，可理解之光之父，我们唤醒与光照之父，使我们受劝回心转意之保证之父。\par

3. 我呼求祢，天主，真理，在祢内、从祢、并通过祢，一切真实的事物都是真实的。天主，智慧，在祢内、从祢、并通过祢，一切智慧的事物都是智慧的。天主，真实而圆满的生命，在祢内、从祢、并通过祢，一切真正而至上活着的事物都活着。天主，真福，在祢内、从祢、并通过祢，一切有福的事物都是有福的。天主，美好与美丽，在祢内、从祢、并通过祢，一切美好而美丽的事物都是美好而美丽的。天主，可理解之光，在祢内、从祢、并通过祢，一切可理解地照耀的事物都可理解地照耀。天主，祢的王国是那感官无法知晓的整个世界。天主，从祢的王国，甚至有一条法律降临到这些低层领域。天主，离开祢就是堕落；回归祢就是复活；居于祢内就是屹立坚定。天主，从祢而出就是死亡；回归祢就是复生；在祢内安居就是活着。天主，除非受骗，无人失落祢；除非被感动，无人寻求祢；除非被洁净，无人找到祢。天主，离弃祢与丧亡是一回事；转向祢与活着是一回事；看见祢与拥有祢是一回事。天主，信德激动我们走向祢，望德举起我们，爱德使我们与祢结合。天主，通过祢我们战胜仇敌，我恳求祢。天主，藉祢的恩赐，我们不致完全灭亡。天主，祢警醒我们守望。天主，祢使我们辨别善恶。天主，祢使我们逃避邪恶，追随良善。天主，通过祢，我们不向灾祸屈服。天主，通过祢，我们忠信地服事并仁慈地治理。天主，通过祢，我们学会那些我们以前视为己有的其实是别人的，而那些我们习惯视为别人所有的其实是我们的。天主，通过祢，邪恶事物的诱惑和勾引无力控制我们。天主，通过祢，减少的财产仍使我们自身完整。天主，通过祢，我们更好的美善不受更差的所辖制。天主，通过祢，死亡被胜利吞灭。天主，祢使我们转向祢。天主，祢脱下我们那不是的，穿上我们那是的。天主，祢使我们堪当被垂听。天主，祢坚强我们。天主，祢引导我们进入一切真理。天主，祢只对我们说善言，既不使我们因恐惧而疯狂，也不允许别人如此行。天主，祢召唤我们回到正途。天主，祢引领我们到生命之门。天主，祢使门向叩门者开启。天主，祢赐给我们生命之粮。天主，通过祢，我们渴慕那饮下后永不口渴的琼浆。天主，祢使世界知罪、知义、知审判。天主，通过祢，我们不被那些拒绝相信的人所困扰。天主，通过祢，我们谴责那些认为在祢面前灵魂没有功过之人的错误。天主，通过祢，我们得以不受软弱和贫乏的元素的束缚。天主，祢洁净我们，并预备我们获得神圣的赏报，求祢慈祥地降临于我。\par

4. 凡是我所说的一切，惟独祢是天主，求祢来帮助我，惟一真实而永恒的本体，那里没有不和、没有混乱、没有变迁、没有匮乏、没有死亡。那里有至高的和谐、至高的确证、至高的稳固、至高的丰盈，以及至高的生命。那里一无所缺，一无所余。那里生者与受生者为一。天主，万物都服事那些服事祢的，每个有德的灵魂都顺从祢。按祢的律法，天穹旋转，星辰运行，太阳赋予白昼生命，月亮调节黑夜：整个宇宙的结构，日复一日因光明与黑暗的交替，月复一月因月亮的盈亏，年复一年因春夏秋冬的依次更迭，周期复周期因太阳轨迹的完成交汇，并经由时间的伟大轨道，不断折叠又展开，正如星辰仍回归它们最初的结合，在这纯然可见的物质所允许的范围内，维持着万物的伟大恒常。天主，按祢永存的律法，变动之事的稳定运动得以不受扰动，循环时代的拥挤进程不断被召回不动宁静的形象：按祢的律法，灵魂的选择是自由的，而善有赏、恶有罚，这些分配是由万物的本性所决定的必然。天主，从祢那里甚至向我们涌来一切益处，由祢一切邪恶被阻止离开我们。天主，在祢之上没有事物，在祢之外没有事物，没有祢就没有事物。天主，万有在祢之下，万有在祢之内，万有与祢同在。祢按祢的肖像和模样造了人，那认识自己的人便发现这肖像和模样。求祢俯听我，俯听我，仁慈地俯听我，我的天主，我的主，我的君王，我的圣父，我的根源，我的希望，我的财富，我的荣耀，我的家，我的祖国，我的健康，我的光明，我的生命。俯听，俯听，仁慈地俯听我，以祢独有的那种方式，这方式虽为少数人所知，但对那少数人却如此熟知。\par

5. 从今以后，惟独祢是我所爱的，惟独祢是我所追随的，惟独祢是我所寻求的，惟独祢是我准备服事的，因为惟独祢是主，以公正的名义，我渴望属于祢的统治。求祢指引，我祈求，并命令祢所愿的一切，但求祢医治并开启我的耳朵，使我听见祢的话语。医治并开启我的眼睛，使我看见祢命令的征兆。从我身上驱散错觉，使我认识祢。告诉我必须趋向何处，才能看见祢，而我盼望，凡祢所吩咐的，我都要去做。主啊，至仁慈的圣父，求祢收留，我祈求，祢的逃犯；我已受够了惩罚，我已长久服事了祢的仇敌（他们在祢脚下），我已长久做了谬误的玩物。求祢收留我从这些中逃离，祢家生的仆人，因为难道不是这些（尽管是另一位主人的）在我逃离祢时接纳了我吗？我感到我必须回归祢：我敲门；愿祢的门为我敞开；教导我通往祢的道路。除意志外我一无所有：除知道那短暂易逝之物应被摒弃、固定永恒之物应被寻求外，我别无所知。圣父，我这样做，因为唯独我知道这些，但我不知道从何处接近祢。求祢教导我，指示我，给我路上的供给。如果是藉着信德，那些投靠祢的人找到祢，那么就赐予信德：如果是藉着德行，就赐予德行：如果是藉着知识，就赐予知识。在我内增加信德、望德与爱德。哦，祢的良善，独一无二且最堪赞叹！\par

7. \Emph{A}. 看，我已向天主祈祷了。\Emph{R}. 那么你想知道什么？\Emph{A}. 我所祈求的一切事。\Emph{R}. 把它们简要地总结一下。\Emph{A}. 天主和灵魂，这就是我所渴望认识的。\Emph{R}. 没有别的了吗？\Emph{A}. 什么都没有了。\Emph{R}. 那么就开始探究吧。但首先请说明，如果天主被呈现给你，你将如何能说：\Quote{够了。}\Emph{A}. 我不知道祂要怎样呈现给我，使我能说\Quote{够了}：因为我相信，我没有任何事物是按我渴望认识天主的样式认识的。\Emph{R}. 那么我们该做什么？你难道不认为，你首先应当知道，什么才是对天主充分的认识，以至于达到那点后，不再进一步寻求？\Emph{A}. 我确实这样认为：但如何实现，我却看不出。因为我曾理解过什么像天主，以致我能说：像我理解这个一样，我也愿意理解天主？\Emph{R}. 既然尚未认识天主，你从哪里知道你不认识任何像天主的事物？\Emph{A}. 因为如果我认识任何像天主的事物，我无疑会爱它：但现在我爱的不外乎天主和灵魂，这两者我都不认识。\Emph{R}. 那么你不爱你的朋友吗？\Emph{A}. 爱他们，我怎能不爱灵魂？\Emph{R}. 那么你同样爱蚊虫和臭虫吗？\Emph{A}. 我说我爱的是赋灵的灵魂，不是动物。\Emph{R}. 那么人要么不是你的朋友，要么你不爱他们。因为每个人都是动物，而你说你不爱动物。\Emph{A}. 人是我的朋友，我爱他们，不是因他们是动物，而是因他们是人，也就是说，因他们被理性灵魂赋予生命，即使在强盗身上我也爱这理性灵魂。因为我完全可以正当地在任何人身上爱理性，即使我正当地恨那滥用我所爱之物的人。因此，我越爱我的朋友，他们越值得地使用他们的理性灵魂，或者当然他们越热切地渴望值得地使用它。\par

8. \Person{理性}：我承认这一点；但若有人向你说：\Quote{我要使你认识天主，如同你认识亚理比乌斯一样}，你岂不感恩并说：\Quote{足够了}吗？\newline{} \Person{奥古斯丁}：我确实会感恩，但我不会说\Quote{足够了}。\newline{} \Person{理性}：请问为什么？\newline{} \Person{奥古斯丁}：因为我对天主的认识甚至不如我对亚理比乌斯的认识，而且我对亚理比乌斯的认识也并不充分。\newline{} \Person{理性}：那么当心，免得你厚颜地渴望在认识天主上得到满足，而你连对亚理比乌斯的认识都尚未满足。\newline{} \Person{奥古斯丁}：\OriginalTerm{非逻辑推论}{Non sequitur}。因为，与星辰相比，我的晚餐何其低贱？然而我不知明天晚餐吃什么；但月亮将处于什么方位，我却可以毫不羞愧地承认我知道。\newline{} \Person{理性}：那么，你认识天主如同你知道明天月亮将在什么方位运行一样，这足够你吗？\newline{} \Person{奥古斯丁}：不够，因为后者我凭感官验证。但我不知道天主或某种隐蔽的自然原因是否会突然改变月亮的常规轨道，若发生，则将使我先前所假定的一切成为虚假。\newline{} \Person{理性}：那你认为这会发生吗？\newline{} \Person{奥古斯丁}：我不认为。但我至少寻求我所能够知道的，而非我所相信的。凡我们所知道的，我们也许可以说我们也相信，但并非凡我们所相信的都知道。\newline{} \Person{理性}：如此说来，你拒绝一切感官的见证？\newline{} \Person{奥古斯丁}：我完全拒绝。\newline{} \Person{理性}：那么你的那位朋友，你说你尚未认识他，你是希望通过感官还是通过理智认识他？\newline{} \Person{奥古斯丁}：凡我通过感官认识他的（若真有什么是通过感官认识的），既微不足道又已足够认识。但那一部分对我怀有情感的部分，即心灵本身，我渴望凭理智认识。\newline{} \Person{理性}：难道还能有别的方式认识吗？\newline{} \Person{奥古斯丁}：绝无可能。\newline{} \Person{理性}：那么你敢称你的朋友，你的至交，为你所不认识的吗？\newline{} \Person{奥古斯丁}：为什么不敢？因为我认为友谊的法律最公平，它规定一个人对他的朋友既不应少于爱，也不应多于爱自己。既然我不认识我自己，那么我宣告他是我所不认识的人，他又何损之有？尤其因为（如我所信）他自己也不认识自己？\newline{} \Person{理性}：既然如此，你所渴望认识的这些事物既然属于理智可得之类，那么当我说你连亚理比乌斯都不认识，却渴望认识天主是厚颜无耻时，你本不应向我提出你晚餐和月亮的类比，因为这些事物如你所说属于感官。\par

9. 但暂且不谈这个，现在回答这点：倘若\Person{柏拉图}与\Person{普罗提诺}关于天主所说的那些话是真的，那么你认识天主如同他们认识祂一样，这样就够了吗？\newline{} \Person{奥古斯丁}：即使承认他们所说的那些话是真的，是否就立刻得出他们知道那些话呢？因为许多人滔滔不绝地讲述他们所不知道的，正如我自己说过，我曾渴望知道我所祈求的一切；倘若我已经知道，我就不再渴望了；但我仍能一一列举它们。因为我所列举的并非我理智所领会的，而是我从各处收集并托付给记忆的，我对之尽其所能地给予充分的信仰；但知道是另一回事。\newline{} \Person{理性}：请告诉我，你至少在几何学中知道什么是线吗？\newline{} \Person{奥古斯丁}：我确实知道这个。\newline{} \Person{理性}：那么你在如此宣称时，难道不害怕学园派吗？\newline{} \Person{奥古斯丁}：毫不害怕。因为他们作为智慧人，不肯冒犯错的风险；但我不是智慧人。因此至今我不羞于宣称我所已经知道的那些知识。但若如我所愿，我终将获得智慧，则我将遵行智慧所示的一切。\newline{} \Person{理性}：我不对此持异议；但我开始要问，正如你知道一条线，你是否也知道一个球，或者说，一个球体？\newline{} \Person{奥古斯丁}：我知道。\newline{} \Person{理性}：两者一样，还是其中一个多，一个少？\newline{} \Person{奥古斯丁}：完全一样。我对两者都完全确定。\newline{} \Person{理性}：你掌握这些是凭感官还是凭理智？\newline{} \Person{奥古斯丁}：不，我曾在此事中把感官当作船来使用。因为当它们把我载到我所瞄准的地方，我便辞退了它们，现在仿佛留在干燥的陆地上，开始在心中思索这些事物，但感官的晃动依然久久在我脑海中游荡。因此我认为在干地上航行，都比凭感官学习几何更容易，虽然年轻的初学者似乎从中得到一些帮助。\newline{} \Person{理性}：那么你不迟疑地将你对这类事物的任何认识称为知识吗？\newline{} \Person{奥古斯丁}：若斯多葛派允许的话，我不反对，因为他们只把知识归于智慧人。我当然坚持我对这些事物有知觉，这是他们甚至承认给愚昧人的；但我对斯多葛派并不十分畏惧：毫无疑问，我以知识持有你问我的那些事物；现在请继续，直到我看到你问我这些的用意。\newline{} \Person{理性}：不要太急切，我们并不赶时间。但请你严格注意，免得你草率让步。我很想给你那种你无所畏惧的事物带来的喜乐，而你却催促赶快，好像无关紧要的事？\newline{} \Person{奥古斯丁}：愿天主成就你预言的结果。所以请随意提问，若我再次犯错，请更严厉地责备我。\par

10. \Emph{理：}那么，显然一条线不可能纵向分成两半？\Emph{奥：}显然不能。\Emph{理：}横切呢？\Emph{奥：}当然可以无限分下去。\Emph{理：}但同样明显的是，如果从一个球体的中心开始任意切分，所得的两个圆永远不会相等？\Emph{奥：}同样明显。\Emph{理：}线和球体是什么？你认为它们是相同的，还是有所区别？\Emph{奥：}谁看不出它们差别很大？\Emph{理：}那么，如果同样清楚地知道二者，而它们又像你所承认的那样彼此迥异，那就必定存在对不同事物的同样知道。\Emph{奥：}谁曾对此有过争议？\Emph{理：}你，刚才就有。因为当我问你，你渴望以何种方式知道天主，以便可以说\Quote{够了}时，你回答说，你无法解释这一点，因为你没有以你所渴望知道天主的那种方式来持有认知，因为你所知道的无一与天主相似。那么又如何？线和球体相似吗？\Emph{奥：}荒谬。\Emph{理：}但我问的不是你知道了什么像天主，而是你知道什么像你渴望知道天主那样。因为你知道一条线，就像你知道一个球体一样，虽然线的属性并非球体的属性。因此请回答，如果你以你知道那个几何球体的方式知道天主，即对天主没有怀疑，就像对那球体没有怀疑一样，这样是否就足够了呢？\par

11. \Emph{奥：}请原谅我，无论你如何激烈地敦促和论辩，我仍不敢说我希望如知道这些东西那样知道天主。因为不仅知道的对象不同，而且知道本身似乎也不相同。首先，因为线和球体并非如此不同，以至于一门科学可以包括对二者的认识；但没有几何学家曾声称教导过天主。其次，如果关于天主的知识与关于这些事物的知识是等同的，那么我知道它们时的喜悦，就应与我确信若我知道了天主时我将会有的喜悦相等。但现在，与祂相比，我对这些东西极度鄙夷，以至于有时我觉得，如果我理解祂，并且是以祂能被看见的那种方式，那么所有这些知识都会从我心中消失；因为即使现在，由于对祂的爱，这些东西几乎不进入我的脑海。\Emph{理：}允许我这样说：你在知道天主时比知道这些东西时会喜悦更多，且多得多，然而并非由于对事物的不同感知；除非我们说，你用不同的视力观看大地和天空的平静，尽管后者的景象比前者更使你感到安慰和愉悦。但除非你的眼睛被欺骗，我相信，如果被问及你是否同样确信看到了大地和天空，你应该回答是，尽管你对大地及其美丽的喜悦远不及对天空的美丽和壮丽。\Emph{奥：}我承认，我被这个比喻打动了，并且被引导承认：正如大地在种类上不同于天空，那些科学论证（尽管真实而确定）也不同于天主那可理解的威严。\par

12. \Emph{理：}你的感动是好的。因为正在与你交谈的理性，承诺会像太阳将自己显示给眼睛一样，将天主显示给你的心灵。因为灵魂的感官就如同心灵的眼睛；而所有科学的确定性，就如同被太阳光照而得以被看见的事物，例如大地及其上的万物；而天主自己就是那光照者。现在，我，理性，在心灵中就犹如观看行为在眼睛中。因为拥有眼睛并不等于在看，看也不等于看见。因此，灵魂需要三样不同的事物：拥有可以善加利用的眼睛、去看、以及看见。健康的眼睛，意味着心灵纯洁，没有肉体的任何污点，即远离并净化了对于必死之物的贪欲。这第一状态，除了信德，没有别的能为她成就。因为对于那被罪恶玷污而疲惫不堪的心灵，还不能展示什么，因为除非健康，她无法看见；如果她不相信否则她就看不见，她就不会关注自己的健康。但倘若她相信情况如我所说，并且如果她要看见，只能在这些条件下看见，但绝望于被治愈；那么她岂不是完全轻视自己，自暴自弃，拒绝遵从医生的嘱咐吗？\Emph{奥：}毫无疑问，尤其是因为疾病中的疗愈必然被视为严厉。\Emph{理：}那么，望德必须加于信德。\Emph{奥：}我如此相信。\Emph{理：}再者，如果她既相信情况如此，又希望自己能被治愈，却不爱、也不渴望那应许的光明本身，反而认为当时应满足于自己的黑暗——这黑暗已由于习惯而变得令她愉悦——她难道不仍然是拒绝医生吗？\Emph{奥：}毫无疑问。\Emph{理：}因此，爱德必须成为第三项。\Emph{奥：}没有什么比这更必要。\Emph{理：}因此，没有这三样，任何心灵都不能被治愈，从而能够看见，即理解它的天主。\par

13. 因此，当心智已然有了健全的眼睛之后，下一步是什么？\Emph{A.} 她要观看。\Emph{R.} 心智观看的行为乃是理性；但因为并非每个观看的人都看见，所以正确而完全的观看，即随之有视力的行为，被称为德行；因为德行即正确或完全的理性。但即使视力本身，即便眼睛已经痊愈，也无力量将双眼转向那光，除非有三样东西常存：信德，灵魂由此相信那件事物——她被要求将目光转向它——具有这样的性质：看见它便赐予幸福；望德，心智由此判断，若她专注观看，便必将看见；爱德，由此她渴望看见并被那景象的享有所充满。专注的观看如今随之带来对天主的真正看见，这是观看的终点；不是因为观看的能力停止了，而是因为它再没有可转向的别处了：这便是真正完全的德行，德行达到了它的终点，随之而来的是幸福的生活。而这看见本身便是灵魂中的那种领会，由领会的主体和被领会的东西所组成：正如用眼看见是由感官与感官对象的结合而产生一样，两者之中任一被撤除，看见便成为不可能。\par

14. 因此，当灵魂得以看见，即领会天主时，让我们看看那三样东西对她是否仍然必要。信德对那现已看见的灵魂何须必要？或望德，当她已然把握时？但爱德非但丝毫不减，反而大大增加。因为灵魂一旦看见了那独特且无伪的美，便会越发爱它，除非她以极大的爱恋将目光固定其上，绝不偏离那景象，否则她将不能停留于那至福的注视中。但当灵魂居于这身体内时，即使她最完全地看见了，即领会了天主，然而，因为肉身的感官仍有其本有的作用，即使它们没有误导的力量，仍不免有某种引发怀疑的能力，因此，那抵抗这些倾向并肯定相反真理的，可被称为信德。此外，在此生中，虽然灵魂已在领会天主中得享福乐，然而，因她忍受许多身体上的烦扰痛苦，她仍有希望：死后这些不便都将止息。因此，望德在此生中也不离弃灵魂。但在此生之后，当她全然将自己聚集于天主内时，爱德依然存留，藉以她被保留在那里。因为她既不受任何假象的干扰，便不能说仍有信德相信那些事物为真；也没有任何事物留待她去盼望，因为她已安全地拥有了一切。因此，有三件事属于灵魂：她得健全，她得观看，她得看见。另外三件事——信德、望德、爱德——对于那三个条件中的前两个，总是必需的；对于第三个，在此生中三者都需要；在此生之后，唯爱德而已。\par

15. 现在请听，就目前所需而论，让我从可感知事物的那个类比出发，也教导一些关于天主的事。即：天主是可理解的，而非可感知的；那些学派的论证也是可理解的，然而它们彼此相差甚远。因为正如地是可见的，光也是可见的；但地若不蒙光照，便不能被看见。因此，那些在学派中所教的事物，凡是理解它们的人丝毫不怀疑它们绝对真实，我们必须相信，除非它们被某种类似其太阳的东西所照亮，否则便不能被理解。因此，正如在这可见的太阳上，我们可观察到三件事：它存在，它照耀，它照明；同样，在那极其遥远而你渴望领会的天主中，也有这三件事：祂存在，祂被领会，祂使其他事物被领会。这两者——天主和你自己——我敢应许我能教你理解。但请回答：你是视这些事为可能的，还是真实的？\Emph{A.} 当然被视为可能的；而且，我必须承认，我曾抱有更大的希望：因为除了那条线和那个球的两种比喻之外，你所说的一切，我没有一件敢说我知道。\Emph{R.} 这不足为奇：因为尚没有一件事被如此陈述，以致要求你领悟。\par

16. 但我们为何迟延？让我们出发：但首先让我们看看（因为这是首要的）我们是否处于健康的状态。\Emph{A.} 请查看，无论是你自己还是我，若你有任何识别应当寻求之事的能力；我将尽我所知回答你所问。\Emph{R.} 除认识天主和你自己之外，你还爱什么别的吗？\Emph{A.} 我本可回答，就我目前的感觉而言，我不爱别的；但更稳妥的说法是我不知道。因为常有这样的事：当我自以为对别的事物无动于衷时，却仍有某种东西闯入我心中，刺伤我出乎意料。又常有些事，思想上想来并不怎么扰乱我，但实际发生时却令我比预想更不安：但现在我自认只有三件事会令我烦恼：害怕失去我所爱的人，害怕痛苦，害怕死亡。\Emph{R.} 因此，你爱的是与至亲相伴的生活、你自身的健康，以及你的肉身生命本身；否则你不会害怕失去这些。\Emph{A.} 确实如此，我承认。\Emph{R.} 那么，如今你所有朋友不全在你身边，你的健康也不很稳固，这便使你心中有些不安。因为我看这是你所说的话中的含义。\Emph{A.} 你看得对；我无法否认。\Emph{R.} 倘若你突然感觉并发现自己身体康健，又看见所有你所爱的人和彼此相爱的人在你陪伴下享受自由安逸，难道你不认为应当适度甚至欣喜若狂吗？\Emph{A.} 适度地，无疑如此。不仅如此，如果这些事如你所说突然临到我，我怎能自制？我甚至怎能掩饰这样的喜悦？\Emph{R.} 由此可见，你仍被心灵的各种病患和动荡所困扰。那么，你竟想以这样的眼睛去看那样的太阳，岂不是厚颜无耻？\Emph{A.} 那么你的结论是：我完全不知道自己健康进展到何种程度，疾病减退了多少，又残留多少。就算我承认这一点。\par

17. \Emph{R.} 难道你没有看见，这肉身的眼睛，即使健康时，也常被这可见太阳的光辉所击，以致不得不转向，退回到自己的昏暗之中？而你现在所打算的，是你被激起要寻求的东西，而不是你渴望看见的东西；然而我正想与你讨论这一点：你认为我们有了多少进步。你没有对财富的欲望吗？\Emph{A.} 至少这已不再是首要的了。因为现在我三十三岁，几乎在过去的十四年间，我已不再渴求财富，即使偶然有财富临到，我也不曾从中寻求什么，只求生活所需和符合自由人身份的使用。西塞罗的一本书彻底说服了我：财富绝不可贪求，但若它们来到我们面前，应以最大的智慧和谨慎来管理。\Emph{R.} 荣誉呢？\Emph{A.} 我承认，只在最近，像是昨天，我才不再渴求这些。\Emph{R.} 妻子呢？你难道不曾被一个美丽、端庄、顺从、有学识，或具有良好的可被你培养的才能的少女的形象所吸引吗？她（因你轻视财富）还会带来一份嫁妆，足以解除你因她而有的负担。此外还暗示你对她不会遭遇不幸。\Emph{A.} 随你如何描绘她，用各种恩赐装饰她，我已决定，没有什么比这样一个伴侣更应避免的了：我看出，没有比女人的媚态和亲昵更摧毁灵魂或身体的男子气概的堡垒了。因此，如果（我尚未发现）繁衍子嗣属于智者的职责，那么任何仅为此目的而进入这种关系的人，在我看来或许值得钦佩，但绝不可效法：因为尝试的危险大于成功的幸福。所以，我相信，我既不违背正义也不违背效用，既不渴求、不寻求、也不娶妻，以此保障我心灵的自由。\Emph{R.} 我现在不问你的决定，而是问你是否仍在挣扎，或已真正克服了欲望本身。因为我们正考虑你眼睛的健康。\Emph{A.} 我丝毫不寻求、不渴望这类事物；我甚至带着恐惧和厌恶回忆它们。你还要什么？这益处与日俱增：因为我越是在盼望看见我所热切渴望的那至高之美中成长，我的爱和喜悦就越是完全转向它。\Emph{R.} 美食呢？你有多在意它们？\Emph{A.} 那些我已决定不吃的，并不诱惑我。至于我尚未戒除的，我承认我在享用它们时感到愉悦，但即便如此，若它们从眼前或口中被撤去，我也不会感到心烦。当它们完全不在时，也没有任何对它们的渴求敢闯入扰乱我的思绪。但无需再问饮食或沐浴：我寻求这些，只求有益于维持健康。\par

18. \Emph{R.} 你有长足的进步；然而，为得见那光而留下的那些东西，仍极大地阻碍着。但我在追求一件在我看来极易表明的事：要么我们毫无需克服之事，要么我们根本毫无进展，而我们所相信已切除的所有那些东西的污染仍存留着。因为我要问你，如果你确信，你能与那群你最亲爱的人一同在智慧的研究和追求中生活，唯一条件是你拥有一份足够充裕的产业，能满足你们共同的一切需要，难道你不会渴望并寻求财富吗？\Emph{A.} 我会。\Emph{R.} 那如果也清楚表明，你将成为许多人的智慧导师，若你的教学权威得到世俗荣誉的支持，甚至你的这些亲人若非自身也享有荣誉，就无法抑制他们的欲望，而这只能通过你的荣誉和尊严才能加给他们呢？那么，荣誉岂不成了值得渴望并努力追求的对象？\Emph{A.} 正如你所说。\Emph{R.} 我不考虑妻子的问题；因为或许不会出现娶妻的必要：尽管如果确定，凭借她丰厚的遗产，所有那些你希望与你安逸地同住一处的人都能得到供养，再加上她由衷的同意，尤其若她出身如此高贵，以至于通过她，那些你刚才在我们的假设中承认必要的荣誉能轻易获得，我不知道你是否有任何责任轻视这些如此得来的好处。\Emph{A.} 但我怎能指望这样的事呢？\par

19. \Emph{R.} 你说得好像我现在正在问你指望什么。我不是在问，被拒绝时不令人快乐的东西，而是问，得到时令人快乐的东西。因为熄灭的瘟疫是一回事，潜伏的瘟疫是另一回事。而且，正如一些智者所说，所有水池都是如此不洁净，以至于它们总是散发一切污秽之物的气味，尽管你并不总是察觉，只有当你搅动它们时才发现。而一个欲望是因无望实现而被压制，还是因灵魂的健康而被驱逐，这两者之间有很大区别。\Emph{A.} 虽然我无法回答你，但你无论如何不能说服我，在我目前所察觉的此种心境中，我毫无所得。\Emph{R.} 这无疑在你看来如此，因为尽管你可能渴望这些东西，但它们在你看来并非因其自身而被渴望，而是为了进一步的目的。\Emph{A.} 这正是我试图说的：因为当我渴望财富时，我是为了这个理由而渴望，即我要富有。而那些荣誉，我已宣称其贪欲刚才被彻底克服，我渴望它们只是出于某种我归于它们的内在光辉的乐趣；当我期待妻子时，除了体面地享受肉欲之外，我别无所求。那时我心中对这些东西有一种真实的渴望；现在我完全鄙视它们：但如果我无法通过这些找到通往我实际渴望之物的道路，我不把它们当作应拥抱之物去追求，而是接受为应容许之物。\Emph{R.} 一个非常出色的区分：因为我也不认为，任何因其他缘故而被寻求的较低级之物的渴望应当受指责。\par

20. 但我要问你，你为什么渴望你所爱之人活着，或者渴望他们与你同住？\Emph{A.} 为了共同和谐地探求天主和我们的灵魂。因为这样，无论谁先发现什么，都能轻易地将同伴引入其中。\Emph{R.} 如果他们不愿探求呢？\Emph{A.} 我会说服他们爱上它。\Emph{R.} 如果你不能，也许因为他们以为自己已经找到，或者认为此类事超乎发现，或者被其他事物的忧虑和欲望所缠绕？\Emph{A.} 我们将尽力而为，我和他们，他们和我。\Emph{R.} 如果连他们的在场也妨碍你的探求呢？你岂不是宁愿选择并努力使他们不与你同住，也不愿在这样的条件下与你同住？\Emph{A.} 我承认正如你所说。\Emph{R.} 那么你渴望他们活着或在场，并非因其本身，而是作为发现智慧的辅助？\Emph{A.} 我完全同意。\Emph{R.} 再者：如果你确信自己的生命阻碍你理解智慧，你还会渴望它延续吗？\Emph{A.} 我会完全回避它。\Emph{R.} 再者：如果你被告知，在这个身体内或离开身体后，你能同样好地获得智慧，你还会在乎是在这一生还是另一生中享受你所极端珍爱之物吗？\Emph{A.} 如果我确知不会经历任何更糟的事，使我从已进步的点上后退，我就不在乎。\Emph{R.} 那么你目前对死亡的恐惧，是基于害怕陷入某种更坏的恶，以致对天主的认识可能从你那里被夺走。\Emph{A.} 我不仅害怕这种可能的丧失——如果我对事实有正确理解的话——也害怕通往我现在急于探索之事的通道被关闭；尽管我相信我已拥有的将留在我身边。\Emph{R.} 因此，不是为了生命本身，而是为了智慧，你渴望生命延续。\Emph{A.} 这是真理。\par

21. \Emph{理}：我们还剩下身体的痛苦，这或许以其自身的力量触动你。\Emph{奥}：我之所以也并非因别的原因而严重畏惧那痛苦，只因它妨碍了我的探究。因为虽然近来我被牙痛的发作折磨得厉害，以致除了我曾致力于学习的事情外，我不能在心里思考任何东西；而当我全部心力都需用于新的进展时，我却完全被阻止去做这些：然而在我看来，如果真理的固有光辉向我显现，我或者就不会感到那痛苦，或者肯定会毫不介意。但是，尽管我从未承受过更严重的事，然而，常常思及我可能必须承受的痛苦有多严重，我有时不得不赞同科尔内略·凯尔苏斯（\Person{Cornelius Celsus}）的说法，即至善是智慧，至恶是身体的痛苦。因为，他说，我们由两部分组成，即灵魂和身体，其中前一部分，即灵魂，是较好的，身体是较差的；因此最高善是较好部分的最好者，最大恶是较差部分的最坏者；现在灵魂中最好的是智慧，身体中最坏的是痛苦。所以，就得出结论，而且我认为是最公正地得出：人的至善是有智慧，至恶是遭受痛苦。\Emph{理}：我们以后会考虑这一点。因为或许智慧本身，我们所追求的智慧，会使我们改变看法。但如果她向我们显示这是真的，那么我们就毫不迟疑地坚持你这个关于至善和至恶的当前判断。\par

22. 现在让我们探讨这个问题：你是怎样一位智慧的爱慕者？你渴望以最纯洁的目光注视她、拥抱她，并以其仅向少数精选的追求者显示的方式，把握她无遮的秀色。因为肯定地，一位激起你炽热爱情的美貌女子，若发现你心中另有所爱，决不会将自己交付于你；那么，那位最纯洁的智慧之美，除非你只为它而燃烧，又岂会向你显现自己？\Emph{奥}：那么为何我仍被置于悲惨之中，被痛苦的渴望所拖延？显然我已经表明，我别无所爱，因为不为自身而被爱的东西并非真爱。如今我至少只为智慧本身而爱她，至于其他事物，如生命、安逸、朋友，我是为了她的缘故而希望它们存在或不存在。但那种美的爱有何限度？在其中我不但不嫉妒别人，反而渴望尽可能多的人与我一同寻求她、注视她、把握她、享受她；因为我知道，我们在所爱之物上结合得越彻底，我们的友谊就越亲密。\par

23. \Emph{理}：确实，智慧所应拥有的正是这样的爱慕者。她所寻求的正是这样的爱慕者，他们的爱除纯洁之外别无他物。但接近她的途径各有不同。因为每人都是按其健康与力量来领悟那唯一而最真实的善。它是某种不可言喻、不可理解的众心灵之光。让这日常之光尽其所能地教导我们有关那更高之光的事。因为有些眼睛如此健康敏锐，以至于它们一睁开，便毫不退缩地转向太阳本身。对这些眼睛来说，光是健康本身，它们不需要教师，或许只需要一个警告。对它们而言，相信、希望、爱就足够了。但另一些人却被那他们急切想看见的光辉所打击，当那景象被撤回时，他们常常愉快地返回黑暗之中。对这些人，虽然他们或许可以合理地被称为健康的，但坚持向他们展示他们尚无力目睹的事物却是危险的。因此，这些人应首先被训练，他们对善的爱要通过延迟来滋养。因为首先应向他们展示那些自身不发光、但可被光照见的东西，如一件衣服、一堵墙之类。然后展示某种自身虽仍不发光，但在光中更灿烂闪耀的东西，如金或银，但尚未耀眼到伤害眼睛。然后或许要谨慎地展示这熟悉的尘世之火，然后是星辰，然后是月亮，然后是逐渐明亮的晨曦，以及发光天空的光辉。在这些事物中，或早或晚，或通过整个序列，或跳过某些步骤，每个人根据自己的习惯训练，最终将毫不畏缩地、满心喜悦地注视太阳。最优秀的导师正是以某种类似的方式对待那些全心全意献身于智慧的人，他们虽然看得模糊，但已经有了能看见的眼睛。因为明智的训练之职责，是以一种逐步接近的方式将人带到她近前，但若不经过这些中间步骤而到达她面前，则是几乎难以置信的福气。但今天，我想我们写得够多了；必须注意健康。\par

24. 又一日来临。\Person{奥}：现在，我恳求你，若你能够，就给我那秩序吧。随你愿引我走何路，经何事，以何方式。把何等艰难、何等辛苦的事加于我，只要它们在我能力之内，我毫不怀疑我必能完成，只要我借此能到达我所渴望之处。\Person{理}：只有一件事我能教导你，别的我不知道。这些感官之事必须完全避开，并且要万分谨慎，免得当我们带着这身体时，我们的翅膀被尘世黏腻的渗液所阻碍，因为我们需它们完整无缺，才能从这黑暗飞向那至高之光。那光甚至不屑于向那些关闭在这身体笼中的人显示自己，除非他们已成为那样的人——无论这笼子是被拆毁还是磨损，他们都会逃入他们本乡的气息中。因此，一旦你变得全然不为任何尘世之事所喜，就在那一刻，请相信我，就在那一刻，你将会看见你所渴望的。\Person{奥}：那将在何时，我恳求你？因为我不认为我能达到这极端的轻蔑，除非我看见了那与这一切相比都是虚无的对象。\par

25. \Person{理}：这样，肉眼可能会说：\Quote{我看见了太阳就不会再爱黑暗。}因为这也似乎属于正确的秩序，实则远非如此。它之所以爱黑暗，是因为它不健康；但除非健康，它无法看见太阳。在这方面，心灵常常犯错：它自以为健康并夸耀自己健康，并且抱怨，好像视力良好，却仍未看见。但那至高之美知道她何时应显现自己。因为她亲自履行医生的职责，比那些被治愈的人更了解谁健康。但我们，就我们所浮现的程度而言，自以为看见了；但我们曾沉溺于何其深的黑暗，或已进步了多少，我们既不允许思考也不允许感受，并且与更深的疾病相比，我们相信自己健康。你难道没看见，昨天我们多么肯定地宣称，我们不再被任何邪恶之事所羁绊，除了智慧之外什么也不爱；并且只为她而寻求或希望别的事物？当我们谈论娶妻的欲望时，那些女人的拥抱对你显得何等卑贱、何等污秽、何等可憎！诚然，在昨夜的同一次交谈中，你感受到那想象的谄媚和那苦涩的甜蜜如何不同于你先前所料，它们搔弄着你；确实远不如往常，但也远非你所想：以至于你那最亲密的医生向你显明了每一件事，既显明你在他的照料下已经走了多远，也显明还有什么有待治愈。\par

26. \Person{奥}：住口，我恳求你，住口。你为什么折磨我？你为什么如此无情地挖掘，并下得如此之深？现在我无法忍受地哭泣，从此我什么都不承诺，什么都不假定；不要再问我这些事。你所说最真实的是：我所渴望看见的那位，祂自己知道我何时健康；让祂按祂所喜悦的行事：祂何时喜悦，让祂显示自己；我现在把自己完全交托给祂的仁慈和眷顾。我一旦坚信，祂绝不会不扶起那些如此倾向祂的人。关于我的健康，我将什么都不说，除非我看见了那美。\Person{理}：确实，不要做别的。但现在止住眼泪，振作精神。你已哭得最厉害，大大加重了你胸中的烦恼。\Person{奥}：你岂要为我的眼泪设定限度，而我看不到我的痛苦有何限度？或者你叫我考虑身体的疾病，而我在内心深处正因令人消瘦的消耗病而憔悴？但是，我恳求你，若你对我有任何能力，试着引我走一些更短的路，好使我至少藉着那光的邻近（若我已有任何进步，我便能忍受）而羞于把眼睛撤回我所离开的黑暗；若我还能说离开了黑暗——那黑暗竟还敢戏弄我的盲目。\par

27. \Person{理}：让我们结束这第一卷吧，若你愿意，以便在第二卷中我们可尝试某种方便呈现的路径。因为你这倾向必须由合理的锻炼来培育。\Person{奥}：我决不容许这卷结束，除非你至少为我打开那光的邻近的一线光芒，我正朝它而去。\Person{理}：你神圣的医生如此满足你的愿望。因为某种光辉攫住了我，邀请我引导你到它那里。因此要专注地接受它。\Person{奥}：引导，我恳求你，把我夺走，任你领向何处。\Person{理}：你确定你渴望认识灵魂和天主？\Person{奥}：那是我全部的心愿。\Person{理}：没有别的？\Person{奥}：完全没有。\Person{理}：什么？难道你不想领悟\OriginalTerm{真理}{Truth}？\Person{奥}：好像不通过她我就能认识这些事似的。\Person{理}：因此她首先须被认识，通过她这些事才能被认识。\Person{奥}：我不拒绝。\Person{理}：那么首先让我们看这一点：既然\OriginalTerm{真理}{Truth}和\OriginalTerm{真实的}{True}是两个词，你认为这两个词所指是两件事，还是一件事？\Person{奥}：我认为是两件事。因为正如\OriginalTerm{贞洁}{Chastity}是一回事，贞洁的（事物）是另一回事，许多事物也是如此；所以我相信\OriginalTerm{真理}{Truth}是一回事，而被宣称为真实的东西是另一回事。\Person{理}：这两者中你认为哪一个更卓越？\Person{奥}：我相信是真理。因为贞洁并非来自贞洁的，而是贞洁的来自贞洁。同样，若有什么是真实的，它确实来自真理而真实。\par

\Emph{R.} 什么？当一个贞洁的人死去，你认为贞洁也死去了吗？\Emph{A.} 绝不。\Emph{R.} 那么，当任何真实的事物消亡时，真理并不消亡。\Emph{A.} 但是，任何真实的事物如何会消亡呢？因为我不明白。\Emph{R.} 我很惊讶你问这个问题：难道我们没有看到成千上万的事物在我们眼前消亡吗？除非你也许认为这棵树要么是一棵树，但不是一棵真实的树，要么如果是真实的，就不能消亡。因为即使你不相信你的感官，并且能够回答说，你完全不知道它是否是一棵树；然而，我相信你不会否认，如果它是一棵树，它就是一棵真实的树：因为这种判断不是来自感官，而是来自理智。因为如果它是一棵假树，它就不是树；但如果它是树，它就必然是一棵真实的树。\Emph{A.} 我承认这一点。\Emph{R.} 那么关于另一个命题：你不承认树是属于那样一种事物，即它产生并消亡吗？\Emph{A.} 我不能否认。\Emph{R.} 因此可以得出结论，某种真实的事物消亡了。\Emph{A.} 我不争辩。\Emph{R.} 接下来呢？难道你不认为，当真实的事物消亡时，真理并不消亡，就像贞洁的人死去时贞洁并未消亡一样吗？\Emph{A.} 我现在也承认这一点，并且急切地等待看你努力要表明什么。\Emph{R.} 那么请注意。\Emph{A.} 我全神贯注。\par

\Emph{R.} 这个命题在你看来似乎是真的吗？凡是存在的，都必然存在于某处？\Emph{A.} 没有比这更令我同意的了。\Emph{R.} 而且你承认真理存在？\Emph{A.} 我承认。\Emph{R.} 那么我们必须探究它在哪里；因为它不在一个地方，除非你也许认为除了物体之外还有别的东西在一个地方，或者认为真理是一个物体。\Emph{A.} 这两件事我都不认为。\Emph{R.} 那么你相信她在哪里？因为她并非无处存在，我们已经承认她存在。\Emph{A.} 如果我知道她在哪里，也许我就不会再寻求什么了。\Emph{R.} 至少你能知道她不在哪里吗？\Emph{A.} 如果你一一审查所有地方，也许我会知道。\Emph{R.} 当然，她不在有死的事物中。因为凡是存在的，都不能存在于某物中，如果那物不在其中存在的话；而真理存在，即使真实的事物消亡，刚才已经承认了。因此，真理不在有死的事物中。但真理存在，又并非无处存在。因此有不死的事物。并且没有真理不在其中的事物是真实的。因此结论是，除了不死的事物之外，没有什么是真实的。并且每一棵假树不是树，假木头不是木头，假银子不是银子，任何虚假的东西，都不是。现在，凡是不真实的，就是虚假的。因此，只有不死的事物才被正确地称为存在。请你仔细考虑这个小小的论证，看其中是否有任何你认为不能承认之点。因为如果它是正确的，我们几乎就完成了我们全部的工作，这在另一卷书中也许会显得更清楚。\par

\Emph{A.} 我非常感谢你，我将勤勉而谨慎地在自己心里审视这些事，而且还要和您一起，在我们安静的时候，如果没有黑暗干扰，并且——我极为害怕——它不会在我心中激起对自身的喜爱。\Emph{R.} 坚定地相信天主，并尽你所能把自己完全交托给祂。不要愿意如同自己拥有自己、自己支配自己；而要承认你是那位最仁慈、最有益的主的奴仆。因为这样祂就不会停止把你提升到祂自己那里，并且不容许任何事临到你，除非那事对你有益，即使你不知道。\Emph{A.} 我听见了，我相信，并尽我所能顺从；并且最恳切地为这件事祈祷，使我获得最大的能力，除非你也许还要求我做什么。\Emph{R.} 暂时这样就好，以后你将会做祂本人——现在已被看见——所要求你的事。\par

\SourceRecord{Translated by C.C. Starbuck.；Nicene and Post-Nicene Fathers, First Series；Vol. 7.；Edited by Philip Schaff.；Buffalo, NY: Christian Literature Publishing Co.,；1888.；<http://www.newadvent.org/fathers/170301.htm>.}

% structure: p0001 source=h1 target=section reason=page-title

\section{《独语录》第二卷}

1. \Emph{A.} 我们的工作已中断了许久，爱是不耐烦的，眼泪也无有止境，除非将所爱的给予爱。所以，让我们开始第二卷吧。\Emph{R.} 让我们开始吧。\Emph{A.} 让我们相信天主将会临在。\Emph{R.} 让我们确实相信，如果这也在我们的能力之内。\Emph{A.} 祂自己的能力就是我们的能力。\Emph{R.} 那么，请尽可能简短而完美地祈祷。\Emph{A.} 永恒不变的天主，让我认识我自己，让我认识你。我已祈祷了。\Emph{R.} 你，将要认识自己的人，你知道你存在吗？\Emph{A.} 我知道。\Emph{R.} 你从何知道？\Emph{A.} 我不知道。\Emph{R.} 你感到自己是单纯的，还是复杂的？\Emph{A.} 我不知道。\Emph{R.} 你知道自己是在运动吗？\Emph{A.} 我不知道。\Emph{R.} 你知道自己在思维吗？\Emph{A.} 我知道。\Emph{R.} 因此，你在思维是真实的。\Emph{A.} 真实。\Emph{R.} 你知道自己是不死的吗？\Emph{A.} 我不知道。\Emph{R.} 在这些你说不知道的事情中，你最渴望知道哪一件？\Emph{A.} 我是否是不死的。\Emph{R.} 那么，你爱活着吗？\Emph{A.} 我承认。\Emph{R.} 当你得知自己是不死的以后，事情将如何？这足够了吗？\Emph{A.} 那确实是一件大事，但对我来说还微小。\Emph{R.} 然而在这微小的事上，你将有多大的喜乐？\Emph{A.} 非常之大。\Emph{R.} 那么，你将不为任何事哭泣了吗？\Emph{A.} 一点也不。\Emph{R.} 如果这一生本身被发现是这样，即它不允许你知道比你已知的更多的东西，你会忍住眼泪吗？\Emph{A.} 不，确实，我会痛哭，以至于生命应该终止。\Emph{R.} 那么，你爱活着不是为了活着的本身，而是为了认识。\Emph{A.} 我承认这个推论。\Emph{R.} 如果这种对事物的认识本身会使你变得不幸呢？\Emph{A.} 我不认为那在任何情况下是可能的。但如果是这样，没有人能是有福的；因为我现在之所以不幸，除了对事物的无知之外，没有其他原因。因此，如果对事物的认识是不幸，那么不幸就是永恒的。\Emph{R.} 现在我看清了你所渴望的一切。因为既然你相信没有人是因为知识而不幸的（由此可知理解力使人有福），但除非活着，没有人是有福的，而除非存在，没有人活着；你希望存在、活着和理解；但存在是为了活着，活着是为了理解。因此，你知道你存在，你知道你活着，你知道你在运用理解力。但是这些事物是否将永远如此，或者这些事物中没有一个将这样，或者某些事物永远存在而某些消失，或者这些事物是否可以减少和增加而一切存在，你渴望知道。\Emph{A.} 正是如此。\Emph{R.} 因此，如果我们证明了我们将永远活着，那么也将推出我们将永远存在。\Emph{A.} 那将推出。\Emph{R.} 那么剩下的就是探讨关于理解力。\par

2. \Emph{A.} 我看到了一个非常清晰简洁的顺序。\Emph{R.} 那么就让这个顺序成为这样：你谨慎而坚定地回答我的问题。\Emph{A.} 我注意听。\Emph{R.} 如果这个世界将永远存在，那么这个世界将永远存在是真实的吗？\Emph{A.} 谁怀疑呢？\Emph{R.} 如果它不会存在呢？那么世界不会存在，难道不是真实的吗？\Emph{A.} 我不争辩。\Emph{R.} 当它毁灭时（如果它将毁灭），难道不是真实的世界已经毁灭了吗？因为只要世界已经终结不是真实的，它就还没有终结：因此，世界已经终结和世界已经终结不是真实的，这是自相矛盾的。\Emph{A.} 这个我也承认。\Emph{R.} 此外，你认为有什么东西可以是真实的，而不是真理吗？\Emph{A.} 绝不。\Emph{R.} 因此，即使万物的框架消失，真理仍将存在。\Emph{A.} 我不能否认。\Emph{R.} 如果真理本身毁灭呢？难道不是真实的真理已经毁灭了吗？\Emph{A.} 这一点谁能否认呢？\Emph{R.} 但若没有真理，真实的东西就不能存在。\Emph{A.} 我刚才已经承认了这一点。\Emph{R.} 因此，真理绝不能消灭。\Emph{A.} 请继续你已开始的讨论，因为没有比这个推论更真实的了。\par

3. \Emph{理} 现在我要你回答我：依你看来，感受和知觉是属于灵魂，还是属于身体？\Emph{奥} 属于灵魂。\Emph{理} 而理智你认为也是属于灵魂的吗？\Emph{奥} 确实如此。\Emph{理} 专属于灵魂，还是也属于别的什么？\Emph{奥} 除了灵魂外，我只看到天主，我相信在祂内存在理智。\Emph{理} 现在我们考虑这一点。如果有人对你说，那墙不是墙，而是一棵树，你会怎么想？\Emph{奥} 要么是他的感官或我的感官出了错，要么是他把墙叫做树了。\Emph{理} 假如他在感官中接受了树的形象，而你接受了墙的形象，难道不能两者都是真的吗？\Emph{奥} 绝不可能，因为同一事物不能既是树又是墙。无论个别事物在我们看来有多么不同，我们之中必然有一人遭受了虚假的想像。\Emph{理} 假如它既不是树也不是墙，而你们两个都错了呢？\Emph{奥} 那确实是可能的。\Emph{理} 因此你刚才忽略了这一点。\Emph{奥} 我承认。\Emph{理} 倘若你承认某事物在你看来与它实际的不同，那么你是在犯错吗？\Emph{奥} 不是。\Emph{理} 因此，显现的可以是虚假的，而看到它的人并不在犯错。\Emph{奥} 可能是这样。\Emph{理} 那么应当承认：看到虚假事物的人并不在犯错，而是同意虚假事物的人犯错。\Emph{奥} 这确实应当承认。\Emph{理} 而这一虚假，何以是虚假的？\Emph{奥} 因为它不同于其所显现的。\Emph{理} 因此，如果没有人在看到它，就没有什么东西是虚假的。\Emph{奥} 推论合理。\Emph{理} 因此虚假不在于事物，而在于感官；但不同意虚假事物的人不会被欺骗。由此推出，我们是一回事，感官是另一回事；因为当感官被误导时，我们可以不被误导。\Emph{奥} 我对此无可反驳。\Emph{理} 但当灵魂被误导时，你敢说你不是虚假的吗？\Emph{奥} 我怎么敢？\Emph{理} 但没有灵魂就没有感官，没有感官就没有虚假。因此，要么是灵魂在运作，要么是与虚假合作。\Emph{奥} 我们之前的推理暗示同意这一点。\par

4. \Emph{理} 现在回答这一点：你认为在将来某个时候，虚假有可能不再存在吗？\Emph{奥} 这对我如何可能？发现真理的困难如此之大，以至于说虚假不存在比说真理不能存在更为荒谬。\Emph{理} 那么你认为，不活着的人能够感知和感觉吗？\Emph{奥} 不可能。\Emph{理} 由此推出，灵魂永远活着。\Emph{奥} 你催我太快进入喜乐了，请慢一点。\Emph{理} 但如果之前的推论正确，我看不出对此事有任何怀疑的理由。\Emph{奥} 我说太快了。因此我更容易相信我曾轻率地做出了某种让步，而不是立即确信灵魂的不朽。尽管如此，请展开这个结论，并说明它是如何得出的。\Emph{理} 你曾说过虚假不能没有感官，而虚假必然存在，因此感官总是存在。但没有灵魂就没有感官，因此灵魂是永存的。而灵魂除非活着，否则无法行使感官。因此灵魂永远活着。\par

5. \Emph{奥} 哦，钝刀！因为如果我曾向你承认，这个世界不能没有人，而这个世界是永恒的，你可能会得出人是永生的结论。\Emph{理} 你观察敏锐。但我们所确立的并非小事，即事物的框架不能没有灵魂，除非在将来某个时候事物的框架中不再有虚假。\Emph{奥} 我确实承认这一推论蕴含其中。但我现在认为，我们应当进一步考虑之前的推论是否在压力下弯曲了。因为我看到，朝着灵魂不朽的方向已经迈出了不小的一步。\Emph{理} 你是否充分考虑了，你是否可能轻率地同意了某事？\Emph{奥} 确实充分考虑了，但我看不出有任何一点可以指责自己轻率。\Emph{理} 因此结论是，事物的框架不能没有活的灵魂。\Emph{奥} 就此而言，只能说明依次有些灵魂出生，有些死亡。\Emph{理} 如果从事物的框架中除去虚假，难道不会导致万物都是真的吗？\Emph{奥} 我承认这一推论。\Emph{理} 告诉我，这堵墙在你看来为什么是真的？\Emph{奥} 因为我没有被它的外表所误导。\Emph{理} 也就是说，因为它正如它所显现的。\Emph{奥} 是的。\Emph{理} 因此，如果某物是假的，是因为它显现得不同于实际；是真的，是因为它正如其所显现；那么，除去看到它的人，就既没有假的，也没有真的。但如果事物的框架中没有虚假，那么万物都是真的。而任何事物只能显现给活的灵魂。因此，如果虚假不能被除去，事物的框架中就有灵魂；如果可以除去，也仍然有灵魂。\Emph{奥} 我看到我们之前的结论确实有所加强，但通过这番扩展我们并未取得进展。因为那个主要动摇我的事实依然存在：灵魂出生和消亡，它们之所以不缺乏于世界，并非由于它们的不朽，而是由于它们的更替。\par

6. R. 你看，有形体的、即可感觉的事物，是否能够被理智所理解？ A. 不能。 R. 那么，天主是否也运用感官来认识事物？ A. 关于此事，我不敢轻率地断言什么；但就猜测所允许的范围而言，天主绝不运用感官。 R. 因此我们断定：灵魂是感官的唯一可能主体。 A. 在或然性许可的范围内，暂且如此断定。 R. 好。那么，你是否承认：这堵墙，如果它不是一堵真正的墙，就不是墙？ A. 再没有比这更让我乐意承认的了。 R. 又是否承认：如果一个形体不是真正的形体，那就不是形体？ A. 这也一样。 R. 因此，如果任何事物若非如其显现即为真，则无物为真；而任何有形之物只能显现于感官之前；而感官的唯一主体是灵魂；且任何形体若非真形体则不能存在；那么，结论就是：若非先有灵魂，便不可能有形体。 A. 你逼我太甚，我已无力抵抗。\par

7. R. 现在请更加留意。 A. 我已准备好。 R. 这当然是一块石头；它之为真，条件在于它不显于外而与其自身相异；而它若不真，就不是石头；而它只能显现于感官之前。 A. 不错。 R. 因此在最隐蔽的地腹之中，或在任何没有感知者存在的地方，就没有石头；这块石头若非我们看见它，也就不是石头；当我们离去，且再无他人看见它时，它也将不再是石头。你即使把箱子锁得很严，无论你在里面藏了多少东西，它们也一无所有。就连木头本身，其内部也不是木头。因为完全不透光的物体深处，逃离一切感官知觉，因此完全没有必然存在。因为如果它存在，就会是真实的；而任何事物若不因如其显现而为真，则非真实；但那个内部不显现；因此它不是真实的——除非你对这个有所反驳。 A. 我看这正是从我先前承认的前提推导出来的；但这结论如此荒谬，以至于我宁愿否定其中任何一条，也不愿承认这是真的。 R. 随你便。那么请想一下，你更愿意说哪一条：有形之物可以显现于感官之外？或者感官可以有灵魂之外的其他主体？或者存在一块石头或其他东西，但它不是真的？或者真理本身需要另外定义？ A. 我求你，我们考虑最后这一条吧。\par

8. R. 那就给\Quote{真}下定义吧。 A. 所谓真，就是事物对认识者而言如其显现，假设认识者愿意且能够认识的话。 R. 那么，任何无人能认识的事物就不是真的？进一步说，如果假就是事物显现异于其实，那么，如果对一个人来说这块石头显现为石头，对另一个人却显现为木头，同一事物岂非既假又真？ A. 前一个说法更让我困扰：如果某事物无法被认识，就会推导出它不是真的。至于同一事物既真又假，我倒不太在意。因为我看同一事物与不同事物相比较，可以既大又小。由此推出，没有任何事物本身是更大或更小的；因为这些是比较的用语。 R. 但如果你说没有事物本身为真，难道你不怕推出没有事物本身存在吗？因为这块木头之所以是木头，正因它是真木头。不可能它本身（即在无认识者的情况下）是木头，却不是真木头。 A. 因此我这样说，这样定义，也不怕我的定义因过于简短而被否决：因为对我而言，真的就是存在的事物。 R. 那么没有东西会是假的，因为任何存在的事物都是真的。 A. 你把我逼入了绝境，我完全无法回答。结果就是，虽然我只愿通过这样的问答来受教，现在却害怕被提问了。\par

9. R. 我们已将自身托付于天主，祂必会施以援手，将我们从这些困境中解救出来，只要我们相信祂，并最虔诚地向祂祈求。 A. 此刻我最乐意做的事莫过于此；因为我从未陷入如此大的黑暗。天主，我们的父，祢劝导我们祈祷，又促成我们向祢祈求；因为我们向祢祈求时，我们活得更好，也更好。求祢在这片幽暗中垂听我这摸索者，向我伸出祢的右手。求祢将祢的光倾注于我，将我召回，不再迷途；求祢亲临我内，好让我也回归于祢。阿们。 R. 现在请尽你所能，以最大的注意力与我同在。 A. 求你说出任何感动你的话，免得我们丧亡。 R. 注意听。 A. 看，我的耳目全都专注于你。\par

10. \Person{理性}：让我们一再地、再三地探讨这个问题：什么是\OriginalTerm{虚假}{falsity}？\Person{奥古斯丁}：我想，除非是那些并不是如其所以为的那样的事物，否则恐怕没有任何东西会被称为虚假。\Person{理性}：更要注意，让我们首先审问感官本身。因为眼睛所见的东西，除非它具有某种真实之物的\OriginalTerm{相似}{similitude}，否则不会被称作虚假。例如，我们在梦中见到一个人，他并非真实的人，而是虚假的，正因为他具有真实之人的相似。因为谁看见一条狗，会有权利说他梦见了人呢？因此，那做梦看见的狗也是虚假的，因为它像一条真狗。\Person{奥古斯丁}：正如你所说。\Person{理性}：此外，如果有人在醒着时看见一匹马，却以为看见了一个人，他不正是被某种人的相似所误导吗？因为如果除了马的形状之外没有别的东西向他显现，他就不能以为看见了人。\Person{奥古斯丁}：我完全承认这一点。\Person{理性}：我们把画中看到的树也称为假树，镜中反射的脸称为假脸，从船上眺望时建筑物的假运动，以及桨插入水中的假折断，所有这些无非是由于其中的\OriginalTerm{逼真}{verisimilitude}。\Person{奥古斯丁}：对。\Person{理性}：同样，我们在双胞胎之间弄错，在鸡蛋之间，在同一枚戒指的印痕之间，以及其他类似的事物。\Person{奥古斯丁}：我全都同意并跟上。\Person{理性}：因此，属于眼睛的事物之相似，是虚假之母。\Person{奥古斯丁}：我不能否认。\par

11. \Person{理性}：但是，除非我弄错，这整个事实的森林可以分为两类。一部分在于同等的事物，另一部分在于低劣的事物。当我们说这一个像那一个正如那一个像这一个时，它们是同等的，比如双胞胎，或者戒指的印痕。当我们说较差者像较优者时，便是低劣的。因为谁照镜子，会愚蠢地说他像那影像，而不是那影像像他呢？这一类又部分在于灵魂所遭受的，部分在于所见的事物。而灵魂所遭受的，要么是在感官中遭受，如建筑物的非真实运动；要么是在自身中从感官所接受的东西遭受，如做梦者的梦，或许还有疯子的幻象。此外，那些出现在我们所看见的事物本身中的东西，有些来自自然，有些是由生物表达和制作的。自然通过\OriginalTerm{生育}{procreation}或\OriginalTerm{反射}{reflection}产生\OriginalTerm{低劣的相似}{inferior similitudes}。通过生育，如父母生出像他们的子女；通过反射，如从各种镜子中。虽然镜子大多是人所制造，但所返回的影像并不是他们制作的。另一方面，生物的作品见于绘画和类似创造物中：其中也可以包括（承认其存在）魔鬼所产生的东西。但是，物体的影子，因为稍加引申可以说它们像其物体且是一种假物体，而且不容否认它们属于眼睛的判断，因此可以合理地归入由自然通过反射产生的那一类。因为每一个暴露在光线下的物体都会反射，并在相反方向投下影子。或者，你看出有什么反对意见吗？\Person{奥古斯丁}：没有。我只是焦急地等待这些例证的结果。\par

12. \Person{理性}：然而，我们必须耐心等待，直到其余感官也向我们报告：虚假寓于真实之相似中。因为听觉中同样有几乎同样多的相似种类：例如，当我们听到一个看不见的说话者的声音时，我们以为是另一个人，因为他的声音相似；在低劣的相似中，\OriginalTerm{厄科}{Echo}是一个见证，或者耳朵本身众所周知的轰鸣，或者时钟中某种模仿画眉或乌鸦的声音，或者做梦者或疯癫者想象自己听到的东西。令人难以置信的是，音乐家所谓的\OriginalTerm{假音}{false tones}有多少为真理作证，这将在后面显示：然而它们（目前只需知道）也并非远离那些人们称为真音的相似。你明白吗？\Person{奥古斯丁}：非常高兴。因为在这里我理解起来毫不费力。\Person{理性}：那么，继续推进：你认为通过嗅觉区分百合与百合容易吗？或者通过味觉区分从百里香采集的蜂蜜（尽管来自不同蜂巢）？或者通过触觉区分鹅绒与天鹅绒的柔软度？\Person{奥古斯丁}：似乎不容易。\Person{理性}：那么当我们梦见自己闻到、尝到或触摸到这类东西时，又是怎样呢？我们不是被一种效果和影像的相似所欺骗，其低劣与空虚成正比吗？\Person{奥古斯丁}：你说得对。\Person{理性}：因此，似乎我们在所有感官中，无论是通过同等的相似还是低劣的相似，要么被欺骗性的相似所误导，要么即使我们没有受骗（如暂停同意，或发现了差异），我们仍将那些我们理解为像真实之物的东西称为虚假。\Person{奥古斯丁}：我不能怀疑。\par

13. \Person{R.} 现在请留意，我们再重述一遍同样的事，好让我们努力要表明的得以更清晰地显现。\Person{A.} 看，我在这里，你尽管说吧。我已决意忍受这迂回的进程，且不会在其中疲倦，因为我热切地盼望最终抵达那我看清我们正趋向之处。\Person{R.} 你做得很好。但请注意，当我们看见蛋的相似处时，我们能否正当地说其中任何一个蛋是假的？\Person{A.} 绝不可能。因为若全都是蛋，它们便是真蛋。\Person{R.} 当我们看见镜中反射的影像时，我们凭哪些标志看出它是假的？\Person{A.} 凭它不可触摸、不发声响、不能独立运动、没有生命，以及无数其他特性——详细列举将很繁琐。\Person{R.} 我看出你厌恶拖延，必须顾及你的急切。那么，不再逐一回忆：假若我们在梦中所见的那些人也能生活、说话、被醒着的人触摸，并且在他们与我们在清醒、健全时所交谈和看见的人之间没有任何差别，我们还有什么理由称它们为假呢？\Person{A.} 我们怎么可能有权利这样做？\Person{R.} 因此，如果它们正是因其与真理最相似而成为真的，并且因为没有任何差别存在于它们与真实和虚假之间（就它们未被那些或其他差别证明为不相似而言），难道不该承认相似是真理之母，不相似是虚假之母吗？\Person{A.} 我无言以对，而且我为自己先前那样仓促的同意感到羞愧。\par

14. \Person{R.} 你若羞愧，可笑了，好像我们选择这种谈话方式不正是为这原因似的：既然我们只与自己交谈，我希望它被称为并题为\WorkTitle{独语录}——这名称虽新，或许刺耳，但用来阐述事实并不不恰当。因为既然真理不能通过更好的方式寻求，只能通过发问与回答，而几乎找不到不以在辩论中落败为耻的人，因此几乎总是发生：当一件事被妥善引入讨论时，却因某些不合理的喧嚷和鲁莽，以及通常隐藏、有时却明显表露的怒气而遭到抛弃；我认为，最为有利且最合乎平安的解决方法是：你决意通过我问你答的方式寻求真理。因此，你若在某个点上不慎束缚了自己，没有理由害怕回头解开绳结；否则便无处可逃。\par

15. \Person{A.} 你说得对；但我完全看不出我错误地同意了哪里：除非是那恰当地称为假的东西与真有某种相似，因为确实没有其他配得上\Quote{假}这一名称的事物出现在我脑中；然而我又被迫承认，那些称为假的事物，正是因其与真不同而被如此称呼。由此得出，那正是那不相似是虚假的原因。因此我心烦意乱；因为我很难记起任何由相反原因产生的事物。\Person{R.} 假若这是宇宙万物中唯一如此的事物呢？难道你不知道，在列举无数动物种类时，唯独鳄鱼在进食时移动上颚？尤其是几乎找不到任何与另一物如此相似，以至于在某个点上又不相像的事物？\Person{A.} 我确实看到这一点；但当我考虑那称为假的东西既与真有相似之处，又有不似之处时，我无法断定它主要在哪一边配得上\Quote{假}的名称。因为若我说：在不相似的那一边；那么就没有什么不能称为假：因为没有什么不与某物不相似，我们承认那是真的。又若我说：应在相似的那一边被称为假；那么不仅那些蛋会因它们惊人的相似而真实地反对我们，甚至我也无法逃脱那种逼迫，即有人强迫我承认万物都是假的，因为我不能否认万物都在某些方面彼此相似。但假使我不怕这样回答：相似与不相似同样赋予称某物为假的权利；你会给我什么出路？因为宣布万物皆假的致命必然性依然会降临到我身上；因为如上所述，万物都发现彼此在某些方面既相似又不相似。我唯一剩下的出路是只称那与所显不同者为假，除非我畏缩于再次遭遇那些我以为早已驶离的怪物。因为一个漩涡又不知不觉地攫住我，把我转回承认那如其所示者便是真。由此得出，没有知者，便没有事物可以是真：在那里，我惧怕在深藏的暗礁上遇难——这些暗礁虽是真实的，却不为人知。或者，若我说那存在者便是真，那么随它谁来反对，结果是没有虚假存在于任何地方。因此我又看见同样的暗礁在前，而且我看出我忍耐你拖延的全部工夫毫无进展。\par

16. \Emph{R.} 请更注意；因为我决不能被说服，我们徒然祈求了神圣援助。因为我看到，尽我们所能尝试了一切之后，我们发现没有什么可留下的，能正当地被称为假的，除了那些要么假装自己不是它本身，要么（包括一切）趋向于存在而并非存在的东西。但前一种虚假要么是欺骗性的，要么是撒谎的。因为那有某种欺骗欲望的东西被正当地称为欺骗性的；这无法被理解为没有灵魂的：但这部分来自理性，部分来自本性；来自理性，在理性受造物中，如在人中；来自本性，在野兽中，如在狐狸中。但我所说的撒谎的，来自那些说出虚假话的人。在这点上他们与欺骗性的人不同，即所有欺骗性的人都试图误导；但并非每一个说出虚假话的人都想误导；因为哑剧、喜剧和许多诗歌充满虚假话，更多是为了愉悦而非误导，而且几乎所有开玩笑的人都说出虚假话。但被正当地称为欺骗性的人，其目的是某人应该被欺骗。但那些不以欺骗为目标，却仍假装一些东西的人，仅仅是撒谎的，或者连这个也不是，至少没有人怀疑他们应被称为愉快伪造者：除非你有些要反对的。\par

17. \Emph{A.} 请继续；因为现在或许你已经开始教导关于虚假，而不是虚假地：但现在我正在考虑你所提到的那类虚假可能是怎样的，即它趋向存在而并非存在。\Emph{R.} 你为什么不该考虑呢？它们就是我们已大量检查过的同一类事物。难道镜中的你的形象不显得想要成为你自己，但因此是假的，因为它不是？\Emph{A.} 这确实似乎如此。\Emph{R.} 至于图画和所有这样被表达出来的相似物，每一件由艺术家制作的东西？它们不都努力成为它们被仿效的那个东西吗？\Emph{A.} 我必须承认这是真的。\Emph{R.} 而且我相信你会同意，做梦者或疯人所受的欺骗属于此类。\Emph{A.} 没有更甚的了：因为没有什么比它们更趋向成为醒着和神志清楚的人所辨认的那些东西；然而它们因此是假的，因为它们所趋向存在的不能存在。\Emph{R.} 现在我何必再说滑动的塔楼、浸入水中的桨、或物体的影子？我认为很明显，它们应以这个规则衡量。\Emph{A.} 它们显然如此。\Emph{R.} 关于其余感官我什么也不说；因为任何人通过思考都不会不发现，在我们感官所觉察的各种事物中，被称为假的是那些趋向成为某物而并非存在的东西。\par

18. \Emph{A.} 你说得对；但我奇怪你为什么要把那些诗歌、玩笑和其他模仿性的琐事从这类中分开。\Emph{R.} 因为的确，愿意成为虚假是一回事，不能成为真实是另一回事。因此，这些人自己的作品，如喜剧、悲剧、哑剧和其他此类东西，我们可以与画家和雕刻家的作品归为一类。因为一个画中的人不能如此真实，无论他多么趋向人的形象，如同那些写在喜剧诗人书中的东西一样。因为他们既不愿意成为虚假，也不是由任何自身的欲望而虚假；而是出于某种必然性，就他们能够追随作者的意图而言。但在舞台上，\Person{罗斯基乌斯}在意愿上是一个虚假的\Person{赫卡柏}，在本性上是一个真实的人；但通过那个意愿，他也是一个真实的悲剧演员，因为他完成了所提议的事：但一个虚假的\Person{普里阿摩斯}，因为他使自己像\Person{普里阿摩斯}，但并不是他。由此现在产生了一件奇妙的事，然而没有人怀疑它是如此。\Emph{A.} 请问，是什么？\Emph{R.} 你怎么看？难道不是所有这些事物在某些方面是真实的，恰恰由于它们在某些方面是虚假的，而且对于它们的真实性质，仅此一点对它们有利，即它们在另一方面是虚假的？因此，它们愿意或应当成为的东西，如果它们避免成为虚假，就绝达不到。因为我提到的那个人，如果他不想成为虚假的\Person{赫克托}、虚假的\Person{安德洛玛刻}、虚假的\Person{海格力斯}和无数的其他东西，又怎能成为真实的悲剧演员呢？或者说，一幅画，例如，怎能是一幅真实的画，除非它是一匹假马？或者，镜中的人的真实形象怎能存在，如果不是一个假的人？因此，如果某事为了成为某种程度的真实，而需要成为某种程度的虚假，那我们为何如此惧怕虚假，并将真理作为至善来追求呢？\Emph{A.} 我不知道，并且我大为惊奇，除非因为在这些例子中，我看不到任何值得模仿的东西。因为为了我们在自己的某种角色上成为真实，我们不应该像演员、镜中映像或\Person{米隆}的铜牛那样被勾画并顺应到扮演另一个人物上；而应当寻求那真理，它并非像被布置在一种两面且自相矛盾的计划上，一面虚假以便另一面真实。\Emph{R.} 你所要求的是崇高和神圣的事情。然而如果我们找到了它们，我们难道不能承认，真理本身由这些事物构成，并且可以说从它们的融合中产生——真理，任何以任何方式被称为真实的事物都从它得名吗？\Emph{A.} 我乐意同意。\par

19. R. 那么你以为如何？论辩术是真的，还是假的？A. 真的，毋庸置疑。但文法也是真的。R. 与前者同一意义吗？A. 我看不出有什么比真更真。R. 那必定是毫无虚假的；有鉴于此，方才你对那些事物感到不悦，那些事物，无论这样或那样，除非是假的，否则不能为真。或者你不知道，所有那些虚构且明显虚假的事物都属于文法？A. 我并非不知道；但是，依我判断，并非通过文法它们才成为假的，而是通过文法，无论它们是什么，都被解释出来。因为戏剧是为了实用或娱乐而编造的虚假。但文法是守护和调节有声言语的科学；其职责包含甚至收集所有已被记忆和文字所托付的人类语言的虚构，并非使它们成为假的，而是教导并强制执行关于这些虚构的某些正确解释的原则。R. 很对。我现在不在乎你是否已经恰当地定义和区分了这些事情；但我要问，是文法本身，还是那种论辩术表明如此。A. 我不否认，定义的力量和技巧——我藉此现在努力区分这些事情——应归于论辩术。\par

20. R. 至于文法本身呢？如果它是真的，它之所以为真，不正是因为它是学问吗？因为\Quote{学问}之名意指某种可学之物；但凡是学习并记住所学的人，不能说不知道；而无人知道虚假。因此，每一种学问和科学都是真的。A. 我看不出同意这个简短的推论有什么冒失之处。但我担心，这会使人认为那些戏剧也是真的；因为这些我们也学习并记住。R. 那么，难道我们的老师不愿意我们相信他所教导的，并知道它吗？A. 不，他完全认真要我们知道它。R. 请问，他是否曾经让我们相信\Person{代达罗斯}飞过？A. 那倒从来没有。但确实，除非我们记住那首诗，他安排得如此周全，以至于我们几乎无法拿住任何东西。R. 那么，你难道否认以下事实为真：存在这样一首诗，并且关于\Person{代达罗斯}有这样的传说流传？A. 我不否认这是真的。R. 那么，你并不否认，当你学习这些时，你学到了真理。因为如果\Person{代达罗斯}飞过为真，而孩童们将此作为虚构的寓言接受并背诵，他们恰恰会因为所背诵的事物是真实的而在心中储存虚假。因为由此导致我们上面所惊异的情况：不可能存在一个关于\Person{代达罗斯}飞行的真实虚构，除非\Person{代达罗斯}飞行是假的。A. 我现在明白了；但我尚未看出这有什么好处。R. 什么好处？除非那个推论不是假的，我们借此得出：除非一门科学是真的，否则它不能成为科学。A. 这有什么意义？R. 因为我想让你告诉我，文法这门科学基于什么：因为科学的真理正基于那个使之成为科学的原则。A. 我不知道如何回答你。R. 你是否觉得，如果其中没有任何定义，没有任何划分和区分为类别和部分，它就绝不可能成为一门真正的科学？A. 我现在明白了你的意思；而且我记不起有任何一门科学，其中定义、划分和推理过程——就它被宣示每个事物是什么而言，无混淆地将其特有属性归于每个类别，既不忽略其特有之物，也不接纳异己之物——不履行那整套功能，科学之名正来源于此。R. 因此，那整套功能，科学正因之称为真。A. 我看出这是必然的。\par

21. R. 现在告诉我，什么科学包含着定义、区分和划分的原则？A. 上面已经说过，这些包含在论辩术的规则中。R. 因此，文法，作为科学，也作为真科学，是由那同一种技艺所创造，这技艺上面已被辩护免于虚假的指控。这一结论无需局限于文法，而可扩展到所有科学。因为你已经说过，而且说得对，你想不起有任何一门科学，其中定义和划分的法则不是其科学特性的基础。但如果它们基于它们之为科学而真，那么谁能否认那使所有科学为真的东西本身就是真理呢？A. 我确实难以不赞同；但这让我犹豫：我们把论辩术本身也算作一门科学。因此，我认为那使这论辩术自身为真的东西才是真理。R. 你的警惕细致非常值得称赞；但我想你不否认，它之为真基于它之为理论和科学。A. 这正是我困惑的原因。因为我注意到它也是一门科学，并因此被称为真。R. 那么？你认为它成为科学，除了其中一切都被定义和划分之外，还有别的根据吗？A. 我无话可说。R. 但如果这种功能属于它，它本身就是和自身就是一门真科学。那么，为何有人会惊讶于那使万物为真的真理是通过它自身并在它自身中为真呢？A. 现在我毫无阻碍地完全赞同那个观点。\par

\Emph{R.} 因此，请注意剩下的几件事。\Emph{A.} 把你所拥有的任何东西都摆出来，只要是我能理解的，我就乐意同意。\Emph{R.} 我们不会忘记，说某物存在于某物之中，有两种含义。一种含义是，它存在于某物中，同时也可以分离而存在于别处，比如这块木头在这个地方，或者太阳在东方。另一种含义是，某物存在于一个主体中，以至于不能与它分离，比如这块木头的形状和可见外观，太阳的光，火的热，心灵中的学问，以及诸如此类。或者您觉得不是这样？\Emph{A.} 这些区分确实是我们最熟悉的，从小就被勤勉地作为思想要素；因此，如果问及这些，我必须立刻承认这个立场。\Emph{R.} 但是，您是否承认，如果主体不持续存在，那么存在于主体中的东西就不能不可分离地持续存在？\Emph{A.} 这一点我也认为是必然的：因为，主体持续存在时，存在于主体中的东西可能不持续存在，任何人稍加思考就能察觉。因为我的身体的颜色，可能由于健康或年龄的原因而发生变化，尽管身体尚未消亡。但这并非对一切事物都同样成立，而是对于某些事物，其与主体的共存并非主体存在所必需的。因为，这堵墙要成为一堵墙，并不必须是它现在呈现的这种颜色；即使偶然变成黑色或白色，或者经历其他颜色变化，它仍然是墙并被称为墙。但是，火如果没有热，那就不是火了；我们也无法说雪是白色的，除非它真的是白色的。\par

\Emph{R.} 那么，我们所寻求的找到了。\Emph{A.} 您是什么意思？\Emph{R.} 你所听到的。\Emph{A.} 那么现在是否已经明确证明，心灵是不朽的？\Emph{R.} 如果您所承认的这些是真实的，那就是以最无可争议的明确性证明了：除非您也许会说，即使心灵死了，它仍然是心灵。\Emph{A.} 至少我绝不会那样说；但我确实要说，正是由于它消亡，它就不再是心灵。我也不会因为这个意见而动摇，因为伟大的哲学家说过，那无论到哪里都赋予生命的东西，不能容纳死亡进入自身。因为，尽管光无论在哪里得以进入，都会使那地方明亮，并且凭借那种令人难忘的对立力量，不能容纳黑暗进入自身；然而它还是会熄灭，那地方因它的熄灭而变暗。因此，那抵抗黑暗的东西，既没有以任何方式接纳黑暗，却通过消亡为黑暗让位，如同它本可以通过离开而让位一样。为此，我担心死亡会以某种方式降临于身体，如同黑暗降临于一个地方；而心灵像光一样，有时离开，有时当场熄灭；因此，现在不是关于身体的每一种死亡都有安全保障，而是必须选择一种特殊的死亡，灵魂通过这种死亡可以被完好地从身体中引导出来，并被指引到一个地方——如果有这样的地方——在那里它不会被熄灭。或者，如果连这也做不到，而心灵如同光一样，是在身体本身中被点燃的，并且没有能力在其他地方持续存在，而每一种死亡都是灵魂或生命在身体中的某种熄灭；那么就必须选择一种死亡方式，以便在允许的范围内，生命在被度过时，可以安全而平静地度过，尽管我不知道如果灵魂死了，这如何可能。哦，那些（无论是出于自身，还是出于任何您愿意的人）已经确信死亡不可怕（即使灵魂消亡）的人，是何等有福！但是，可怜的我，至今没有任何推理、任何书籍能够说服我这一点。\par

\Emph{R.} 不要叹息，人的心灵是不朽的。\Emph{A.} 您如何证明？\Emph{R.} 从您上面非常谨慎地承认的那些事。\Emph{A.} 我确实不记得在回答您的提问时，我的承认有任何疏忽之处；但现在，我请求您，把所有事汇总成一句话；让我们看看经过这么多迂回之后，我们到了哪一步，而且我不希望您在此过程中再问我问题。因为，如果您要简明扼要地列举我所承认的那些事，为什么还需要我的回应？或者，如果实际上我们取得了任何可靠的结果，您故意用延迟的喜悦来折磨我吗？\Emph{R.} 我将会做我看出您所希望的事，但请最专心地注意。\Emph{A.} 现在说吧，我在这里；您为什么杀我？\Emph{R.} 如果每一件存在于主体中的东西都始终持续存在，那么主体本身必然也始终持续存在。而每一门学问都存在于主体心灵中。因此，如果学问永远持续，那么心灵必然永远持续。现在，学问就是真理，而正如本书开头理性曾说服您的，真理始终存在。因此，心灵永远持续，而且死了就不能被称为心灵。因此，唯有能够证明前面任何承认是没有理由的人，才能避免在否认心灵不朽时陷入荒谬。\par

25. \Emph{A.} 现在我已准备好投入那预期的喜乐，但有两个思绪使我犹豫不决。首先，我们绕了这么大一个圈子，遵循着我不知道什么样的推理链条，而整个谈论的话题本可以如此简短地证明，正如刚才所展示的那样，这让我感到不安。因此，我担心这谈论如此漫长地采取如此谨慎的步伐，仿佛有意设置埋伏。其次，我不明白一门学问如何总是存在于心智中，因为一方面，熟悉它的人如此之少；另一方面，凡是知道它的人，在早期的童年时代也有那么长的时间对它一无所知。因为我们既不能说未受教育者的心智不是心智，也不能说那门他们一无所知的学问存在于他们的心智中。如果这完全荒谬，那么结果要么是这门学问并不总是存在于心智中，要么是那门学问并非真理。\par

26. \Emph{R.} 你可以注意到，我们的推理绕了这么大一个圈子并非徒劳。因为我们一直在探究什么是真理，即使在现在，在这个思想与事物的森林中，它引诱我们的脚步进入无数条途径，我们也未能，如我所见，最终追踪到它。但我们要做什么呢？我们是否应放弃我们的尝试，等待希望某本书落入我们手中，以满足这个问题？因为我认为，许多人在我们这个时代之前已经写过，我们并未读过；而现在，为了对我们不知道的事不做猜测，我们清楚地看到，关于这个主题有很多著作，既有诗体的，也有散文的；并且出自那些作品不可能为我们所不知的人，他们的才华我们知道是如此，以至于我们不能绝望于在他们的作品中找到我们所需的：尤其是当我们眼前就有此人，我们认出了我们曾哀悼为已死的雄辩在他身上已重新焕发活力。他会在著作中教导我们真正的生活方式之后，还让我们对生活的真正本质保持无知吗？\Emph{A.} 我确实不这么认为，并对此抱有很大希望，但我有一件悲伤的事，就是我们没有机会向他敞开我们对他的热忱，或对智慧的热忱。因为他肯定会怜悯我们的渴望，并比现在更快地倾注。因为他是安稳的，因为他现在已确信灵魂的不朽，也许不知道还有那些人，他们太过深切地经历了这种无知的不幸，而不援助他们是残酷的，尤其是当他们恳求时。但另一位确实从长久的熟悉中知道我们的渴望热情；但他离得如此之远，我们的处境如此，以至于我们几乎连寄一封信给他的机会都没有。我相信他最近在阿尔卑斯山外的退隐中编了一道咒语，在其禁令下，死亡的恐惧被迫逃离，被持久冰封而变得麻木的灵魂冷僵被驱散。但与此同时，当这些帮助缓慢地运抵此处——这是一种我们无法命令的恩惠——我们的闲暇正在浪费，我们的心智完全依赖于他人意志的不确定决定，这不是最不值得的吗？\par

27. 对此我们该说什么呢？我们曾恳求天主，并仍在恳求，求祂为我们指示一条道路，不是通往财富，不是通往肉体的快乐，不是通往世俗的荣誉和官位，而是通往对我们自身灵魂的认识，并同样向寻求祂的人启示祂自己。祂果真会抛弃我们吗？或者我们该抛弃祂吗？\Emph{R.} 祂确实最不可能抛弃那些渴望这类事物的人；因此，我们若抛弃如此伟大的向导，也应该是我们思想中奇怪的事。所以，如果你愿意，让我们简要回顾一下那些考虑因素，从中得出两个命题之一：要么真理永远长存，要么真理是论证的理论。因为你说过这些观点在你的心中摇摆，使我们整个事情的最终结论不那么确定。或者我们不如探究这个问题：一门学问如何能存在于一个未受训练的心智中，而我们又不能否认它是一个心智？因为这一点似乎让你感到不安，以至于使你对先前的让步再次产生怀疑。\Emph{A.} 不，让我们先讨论前两个命题，然后我们再考虑这后一个事实的性质。因为如此，依我判断，就不会有余留的争议了。\Emph{R.} 就这样吧，但要极其专注和谨慎地注意。因为我知道你倾听时会发生什么，那就是当你过于专注于结论，并期望它现在、现在就要得出时，你对我问题中隐含的要点未经充分仔细的审查就同意了。\Emph{A.} 也许你说的是真话；但我会尽我所能对抗这种疾病：现在就请开始问我吧，免得我们在多余的事情上耽搁。\par

28. \Emph{R.} 我记得，从这个真理——真理不能消亡——我们得出结论，不仅如果整个世界消亡，甚至如果真理本身消亡，世界和真理都已消亡这一点仍会是真实的。现在，没有真理就没有真实的东西：因此真理绝不会消亡。\Emph{A.} 我承认这一切，如果它被证明是假的，我会非常惊讶。\Emph{R.} 那么让我们考虑另一点。\Emph{A.} 请允许我稍微思考一下，免得我很快又困惑地回来。\Emph{R.} 那么，真理已经消亡这件事不会是真实的吗？如果它不是真实的，那么真理就没有消亡。如果它是真实的，那么在真理陨落之后，真实的东西在哪里？既然现在没有真理了。\Emph{A.} 我不再需要思考和考虑；继续谈其他事吧。确实，我们会尽力安排，使博学和智慧的人能读到这些沉思，如果他们发现任何轻率之处，可以纠正它们：因为对我自己而言，我不相信现在或将来我能发现有什么可以反驳这一点。\par

29. \Person{R.} 那么，真理之所以被称为真理，难道不是因为它使一切真实的东西成为真实吗？\Person{A.} 没有别的原因。\Person{R.} 它被称为真实的，除了因为它不是虚假的，还有别的理由吗？\Person{A.} 怀疑这一点就是疯狂。\Person{R.} 那与任何事物的相似相符，却又不是它所显现的那事物的相似，难道不是虚假的吗？\Person{A.} 我确实没有看到任何我更愿意称之为虚假的东西。然而，通常被称为虚假的，是与真实的相似相距甚远的东西。\Person{R.} 谁否认这一点？但即便如此，因为它包含着对真实的某种模仿。\Person{A.} 怎么讲？因为当人们说美狄亚驾着有翼的蛇车飞走时，那件事丝毫没有模仿真理；因为它根本不存在，而且当它本身绝对虚无时，也无法模仿任何东西。\Person{R.} 你说得对；但你没有注意到，那绝对虚无的东西甚至不能被称为虚假。因为如果它是虚假的，它就是存在的；如果它不存在，它就不是虚假的。\Person{A.} 那么，我们难道不能说那个关于美狄亚的怪诞故事是虚假的吗？\Person{R.} 当然不能；因为如果它是虚假的，它怎么会是一个怪诞的故事呢？\Person{A.} 妙极了！那么，当我说\par

\Quote{我将巨大的有翼之蛇套在我的车上，}\par

难道我说的不是虚假的吗？\Person{R.} 你确实说了，因为那就是你所说的虚假的东西。\Person{A.} 请问，是什么呢？\Person{R.} 就是那诗句本身所包含的句子。\Person{A.} 请问那又有什么对真理的模仿呢？\Person{R.} 因为即使美狄亚真的做了那件事，这诗句也会具有同样的内容。因此，就其措辞而言，一个虚假的句子模仿了真实的句子。如果它不被相信，它就在这一点上模仿了真实的句子：它被表达得像真实的句子一样，并且它只是虚假的，而不是误导性的。但如果它获得了信任，它还模仿了那些因为是真实的而被相信的句子。\Person{A.} 现在我明白了，我们所说的事物与我们谈论某事的事物之间有很大的区别；因此我现在表示赞同：因为只有一个命题阻止了我，那就是我们称之为虚假的东西，除非它包含对某事物的模仿，否则不能正确地称为虚假。因为，谁若称一块石头为假银，难道不该被嘲笑吗？然而，如果有人声称一块石头是银子，我们就说他说话虚假，也就是说，他说了一个虚假的句子。但我认为，我们称锡或铅为假银并非不合理，因为那事物本身仿佛模仿了银子；并且我们对此所做的陈述因此并非虚假，而是那被陈述的事物本身是虚假的。\par

30. \Person{R.} 你理解得很好。但考虑这一点，我们能否也合宜地称银为假铅？\Person{A.} 依我之见，不能。\Person{R.} 为什么？\Person{A.} 我不知道；只是我觉得那样称呼完全违背我的意愿。\Person{R.} 或许是因为银是更好的东西，这样的称呼会轻蔑它；但若称铅为假银，却似乎赋予铅某种荣誉？\Person{A.} 你恰如其分地表达了我所想。因此我相信，那些穿着女装招摇过市的人被定为臭名昭著并失去作证资格，是很有道理的；我不知道更合宜地称他们为假女人，还是假男人。然而，真正的演员，以及真正臭名昭著的人，无疑我们可以这样称呼他们；或者，如果他们隐藏不露，而臭名昭著意味着恶名，我们也可以不无真实地称他们为真正的无用之物。\Person{R.} 我们会有别的机会讨论这些事情：因为有许多事，仅就其外表看来是卑下的，然而为了某个可称赞的目的而做时，却显得高尚。这是一个大问题：一个人为了解救他的国家，是否应该穿上女人的衣服去欺骗敌人，而正是由于他是一个假女人，反而可能显明为更真实的男人；以及一个明智的人，若以某种方式确实确知他的生命将对人类利益是必要的，是否应该宁愿冻死，也不穿上女装（如果没有别的方法）。但关于这一点，如上所述，我们以后会考虑。因为毫无疑问，你看出它需要多么仔细的探究，才能在不陷入各种不可原谅的卑劣行径的前提下，做到何种程度。但现在——这足以解决眼前的问题——我认为现在很明显，而且毫无疑问，没有什么是虚假的，除非通过对真实的某种模仿。\par

31. \Emph{A.} 继续讨论余下部分；因为我对此深信不疑。\Emph{R.} 那么我问这个问题：除了我们所受教的科学（其中智慧本身的研究也应被包括在内）之外，我们能否找到某种如此真实的东西，以至于它不像舞台上的\Person{阿基里斯}那样，在一方面是假的，以便在另一方面可以是真的？\Emph{A.} 对我来说，确实似乎可以找到许多这样的东西。因为没有任何科学包含这块石头，而且，为了使它是一块真的石头，它也不模仿任何东西从而被称为假的。提到这一点，你看就有机会详述无数事物，它们自己出现在思想中。\Emph{R.} 我明白，我明白。但是它们在你看来是否可以被包含在身体的单一名称之下？\Emph{A.} 它们可能会这样显得，如果我或者确知\OriginalTerm{虚空}{Inane}是虚无，或者认为心灵本身应被算在身体之中，或者相信天主也是一个身体。如果所有这些都是，我看它们并不通过模仿任何东西而成为假的或真的。\Emph{R.} 你让我们踏上长途，但我会以最简洁的速度进行。因为确实你所谓的虚空是一回事，你所谓的\OriginalTerm{真理}{Truth}是另一回事。\Emph{A.} 确实天差地别。因为如果我认为真理是某种虚空，或者如此热切地寻求某种虚空，那还有谁比我更虚空？因为我所渴望找到的除了真理还有什么？\Emph{R.} 因此或许你也承认这一点：没有任何东西是真实的，除非它因真理而成为真实。\Emph{A.} 这在此前阶段就已经显明了。\Emph{R.} 你是否怀疑除了虚空本身之外没有任何东西是虚空的，或者当然身体不是虚空？\Emph{A.} 我毫不怀疑。\Emph{R.} 因此我推测，你相信真理是某种身体。\Emph{A.} 绝无可能。\Emph{R.} 什么是身体？\Emph{A.} 我不知道；没关系：因为我认为你知道，即使是那个虚空，如果它是虚空的话，在没有身体的地方更加完全地如此。\Emph{R.} 这当然是明显的。\Emph{A.} 那我们为什么还拖延？\Emph{R.} 那么你看来是真理创造了虚空，或者任何没有真理的地方有什么真实的东西吗？\Emph{A.} 两者似乎都不真实。\Emph{R.} 因此虚空不是真实的，因为它既不能由非虚空的东西变成虚空；而且显然，缺乏真理的东西不是真实的；最后，那个被称为虚空的东西之所以如此称呼，是因为它是虚无。那么，那不存有的东西怎能是真实的？或者那绝对虚无的东西怎能存有？\Emph{A.} 好吧，那我们就将虚空当作虚空而抛弃吧。\par

32. \Emph{R.} 关于其余部分，你怎么说？\Emph{A.} 什么？\Emph{R.} 因为你看我这一方有多少依据。因为我们剩下\OriginalTerm{灵魂}{Soul}和\OriginalTerm{天主}{God}。如果这两个是真实的，因为真理在它们里面——关于天主的永生，无人怀疑。但心灵被认为不朽，如果那不灭的真理被证明在它里面。因此让我们考虑这最后一点：身体是否不是真正真实的，即，在它里面是否有，不是真理，而是某种真理的影像。因为即使在我们已知会朽坏的身体中，我们找到像科学中那样的真理元素，那么讨论技艺是真理（一切科学借以成为真实）的推论就不那么确定了。因为连身体也是真实的，这身体似乎不是由论证的力量形成的。但是，如果连身体也是通过某种模仿而真实，并因此不是绝对纯粹真实的，那么，也许就没有什么能阻碍论证理论被教导为真理本身。\Emph{A.} 与此同时，让我们考察一下身体；因为即使这个问题解决了，我也看不到结束这场争论的前景。\Emph{R.} 你从哪里知道天主的旨意？因此请听：因为至少我认为身体被包含在某种形式和外形之中，如果它没有这些，它就不是身体；如果它真实地拥有这些，它就会是心灵。还是事实并非如此？\Emph{A.} 我部分同意，其余部分我怀疑；因为，除非某种形状被保持，我同意它不是一个身体。但是，如果它真实地拥有形状，它就会是心灵，这个我不太理解。\Emph{R.} 那么你不记得关于本书绪论和你的那个几何学的任何事情了吗？\Emph{A.} 你提得很及时；我确实记得，而且非常愿意记起。\Emph{R.} 在身体中能找到像那门科学所演示的那样的图形吗？\Emph{A.} 不，它们被证实有多大的劣等，简直难以置信。\Emph{R.} 那么，你认为哪个是真实的？\Emph{A.} 我请求你不要认为有必要甚至对我提出这个问题。因为有谁如此迟钝，以至于看不出在几何学中教导的那些图形居于真理本身之中，甚至真理居于它们之中；而那些有形的图形，由于它们似乎趋向于这些图形，有我不知道的真理的模仿，因此是虚假的？因为现在你努力要展示的那整个事理，我明白了。\par

33. \Emph{R.} 除了探讨论辩的学问之外，还有什么更长的必要呢？因为无论几何图形是在真理之中，还是真理在它们之中，它们都包含在我们的灵魂中，也就是在我们的理智中，这一点没有人质疑；而且通过这一事实，真理也必然存在于我们的心智中。但如果每一种学问都如此存在于心智中，如同不可分离地存在于主体中，并且如果真理不能消亡，那么，我问，为什么我们因着某种与死亡的亲昵而对心智的永生抱有怀疑呢？或者，那线条、正方或圆形是否还有其他东西，它们模仿这些以便成为真实的？\Emph{A.} 我绝不能相信这一点，除非或许线条另有所指，即无宽度的长度，而圆则是从中心到各边等距的封闭线条。那么我们还犹豫什么？难道真理不在这些事物所在之处吗？\Emph{A.} 愿天主远离这种疯狂。\Emph{R.} 难道学问不在心智中吗？\Emph{A.} 谁会这样说呢？\Emph{R.} 但是，如果主体消亡，那存在于主体之中的东西能够持存吗？\Emph{A.} 我何时能想象这样的事呢？\Emph{R.} 剩下的只能是假设真理会失败。\Emph{A.} 这怎么可能发生呢？\Emph{R.} 因此，灵魂是不朽的：现在终于向你的论证屈服吧，相信真理；她大声宣告她居住在你内，并且是不朽的，她的居所不会因身体的任何死亡而被夺走。远离你的阴影，返回你自身；你所恐惧的毁灭毫无意义，除非你忘记了你是不可能被毁灭的。\Emph{A.} 我听到了，我开始恢复理智，开始回想自己。但我恳求你加快处理剩余的问题：在一个未受训练的心智中——我们不能称它为有死的——学问和真理如何被理解。\Emph{R.} 这个问题需要另一卷书，如果你想彻底探讨它；而且，我还看到你有机会回顾我们已经尽最大努力探讨过的那些问题；因为如果所承认的每一个问题都无可置疑，那么我认为我们已经取得了很大成就，并且可以相当有把握地继续深入探究。\par

34. \Emph{A.} 正如你所说，我心甘情愿地服从你的指示。但有一点我恳求你，在你给本书定下结束之前，简明扼要地说明一下，那包含在理智中的真实图形与思想为自己构想的图形（希腊文称为\OriginalTerm{幻想}{Phantasia}或\OriginalTerm{幻象}{Phantasma}）之间的区别是什么。\Emph{R.} 你所寻求的，是只有最纯洁视力的人才能看见的东西，而你对这件事的视力训练不足；现在我们在这广阔的范围内别无他图，只是为了训练你，使你能够看见；然而，如何能教导说区别非常大，也许我可以稍加努力使之清楚。因为假设你忘记了某事，而别人希望你回忆起来。他们于是说：是这个吗，是那个吗？提出了与它不同却看似相似的东西。但你既没有看到你想回忆的东西，却又看到这不是所提出的东西。当这种情况发生时，你是否认为它等同于完全的遗忘？因为这种区分的能力，借此，对错误提出的建议予以拒绝，正是回忆的一部分。\Emph{A.} 似乎如此。\Emph{R.} 因此，这样的人尚未看见真理，但他们不会被误导和欺骗；他们足够知道他们所寻求的。但如果有人说，你出生后几天你就笑了，你不敢说那是假的；如果他是可信的权威，你准备不是记住，而是相信；因为那整段时间对你来说是埋葬在最真实的遗忘中。或者你有不同看法？\Emph{A.} 我完全同意这一点。\Emph{R.} 因此，这种遗忘与那种遗忘截然不同，但那是一种中间状态。因为还有另一种更接近、更邻近真理的回忆和重新点燃的视觉：如同我们看见某物，并确信我们曾在某个时候见过它，并且断言我们知道它；但是，在哪里、何时、如何、与谁一起认识了它，我们却费尽心思在记忆中寻找答案。如同发生在一个人身上，我们也问我们在哪里认识了他：当他记起时，突然整个事情像一道光一样闪现在记忆中，我们不再费力去回忆。这种遗忘你不知道吗，或者不明显吗？\Emph{A.} 还有什么比这更明白、或者更频繁地发生在我身上呢？\par

35. \Emph{理性}凡在博雅技艺中受过良好训练者，就是这样；他们借着学习，将埋没于自身遗忘中的技艺发掘出来，仿佛重新挖出；他们不满足，也不停息，直到他们最广泛、最清晰地看见真理的全貌——在那些技艺中，已有某种光辉向他们闪耀。但由此，某些虚假的色彩与形态仿佛涌入思维之镜，常误导探究者，欺骗那些以为他们所知或所探究的便是全部的人。这些想象本身须极其小心地避免；因其随思维之镜的变化而变化，故可被察觉为虚妄，而真理之面则恒一不变。例如，思维为自己描绘一个或大或小的方形，并将其呈现在眼前；但渴望看见真理的内心，若能如此，便转而趋向那普遍概念，依此判断所有这些都是方形。\Emph{奥古斯丁}若有人对我们说，心灵是根据眼睛所见的习惯来判断，那又如何？\Emph{理性}那么，它为何判断（假设它受过良好训练）一个任何大小的真正球体可以被一个真正平面切于一点？眼睛何曾见过，或如何能见到这样的事？因为这类事物无法在纯粹的思想想象中成形？当我们甚至在心中描画最小的想象中的圆，并从它向中心画线时，难道我们不是证明了这一点吗？因为当我们画了两条线，其间几乎容不下一根针尖时，我们即使在想象中也不能再在它们之间画其他线，使它们无任何混合而到达中心；而理性却宣称可以画出无数条线，彼此仅在中心相切，以至于它们之间的每一间隔甚至可画一个圆。既然幻象不能做到这一点，且比眼睛本身更为欠缺（因为正是通过眼睛，它才加于心灵），那么它显然与真理大不相同，并且当幻象被看见时，真理并不被看见。\par

36. 当我们开始讨论理智的能力时——这主题是我们所提出要加以阐述和探讨的，当有任何关于灵魂生命的事令人担忧时——将更费力、更精细地处理这些要点。因为我相信你并非不惧怕人的死亡，即使它不杀死灵魂，也可能导致对万物的遗忘，以及对真理本身的遗忘——假如曾有真理被发现的话。\Emph{奥古斯丁}这邪恶有多么可怕，无法言表。因为永生将是怎样的呢？或者，如果灵魂像我们在初生婴孩身上所见的那样活着，那什么样的死亡不是比它更可取呢？更不用说在母腹中所活的生命了；因为我并不认为它不存在。\Emph{理性}鼓起勇气；现在我们所感受到的天主，必与寻求祂的人同在，祂应许在此世之后赐予一个极其荣福之身，以及毫无虚妄的真理的圆满。\Emph{奥古斯丁}但愿如我们所望。\par

\SourceRecord{Translated by C.C. Starbuck.；Nicene and Post-Nicene Fathers, First Series；Vol. 7.；Edited by Philip Schaff.；Buffalo, NY: Christian Literature Publishing Co.,；1888.；<http://www.newadvent.org/fathers/170302.htm>.}
